% !TEX program = xelatex
% 编译方式：xelatex -> bibtex -> xelatex -> xelatex
% 参考文献数据库与综述共用：vla_fpga_survey.bib（同目录）
\documentclass[a4paper,10.5pt]{ctexart}

\usepackage[margin=2.4cm]{geometry}
\usepackage{amsmath,amssymb}
\usepackage{booktabs,tabularx,multirow,longtable,array}
\usepackage{graphicx}
\usepackage{xcolor}
\usepackage{caption}
\usepackage{enumitem}
\usepackage{tikz}
\usetikzlibrary{arrows.meta,positioning,fit,shapes.geometric,calc}
\usepackage{pgfplots}
\pgfplotsset{compat=1.17}
\usepackage[numbers,sort&compress]{natbib}
\usepackage[colorlinks=true,linkcolor=blue!50!black,citecolor=green!40!black,urlcolor=blue!60!black]{hyperref}

\captionsetup{font=small,labelfont=bf}
\setlist{nosep,leftmargin=2em}
\renewcommand{\arraystretch}{1.15}
\newcolumntype{Y}{>{\raggedright\arraybackslash}X}

% ---------- 数据来源标注（与综述一致） ----------
\newcommand{\tM}{\textsf{[M]}}              % 简报：板上实测或周期精确 RTL 仿真
\newcommand{\tP}{\textsf{[P]}}              % 简报：模型预测
\newcommand{\tMsim}{\textsf{[M, sim]}}      % 简报：RTL 仿真实测
\newcommand{\tMsw}{\textsf{[M, software]}}  % 简报：软件模拟实测
\newcommand{\tD}{\textsf{[派生]}}           % 由简报数字直接算术得到
\newcommand{\tE}{\textsf{[估算]}}           % 依据公开模型结构所作量级估计，不是项目结果
\newcommand{\tR}{\textsf{[建议]}}           % 本报告提出的设计或验证建议，尚未实施
% 待上板测量标记：编号对应第 \ref{sec:board} 节“待上板测量清单”
\newcommand{\tB}[1]{\textcolor{red!70!black}{\textsf{\textbf{[待上板测 #1]}}}}
\newcommand{\pizero}{$\pi_0$}

\title{\textbf{\pizero{} 视觉-语言-动作模型单片 FPGA 实现：\\设计方案与验证分析报告}}
\author{（作者信息待填）}
\date{}

\begin{document}
\maketitle

\begin{abstract}
\noindent
本报告面向 FPGA 与体系结构技术评审，系统说明一个在单片 Achronix Speedster7t AC7t1500 上端到端执行 \pizero{} 视觉-语言-动作（VLA）模型整次推理（一个动作块，下称 chunk）的设计方案与验证体系，并给出演示（demo）方案。报告首先从需求与负载出发，推导出“多节点自主阵列 + 列并行 MLP72 级联链 + 独立向量节点 + 编译器生成程序”的总体架构；随后逐层说明 GEMM 链、PV 链与向量节点的微结构，W8A8 数值方案及其位真模型，chunk 生成器的编译与调度流程，以及在 Speedster7t 上的物理实现、时钟域划分与布线约束。验证部分建立了从 fp32 参考、量化位真模型、编译器期望内存、周期精确 RTL 仿真到硅片逐拍核对的五级验证链，并给出各级的覆盖目标与判据。在最佳 INT8 设计点上，完整 35 节点阵列的整 chunk 时延为 1.221~s，46 帧录制数据上关节角最大误差 $0.665^\circ$、平均 $0.072^\circ$，结果与编译器期望内存逐位一致；同一 chunk 在 NVIDIA Jetson AGX Thor 上以 bf16 编译为 125.3~ms。按校准模型，向量工作与串行权重装载合计约占时延的 65\%，布线而非乘法器是阵列规模的约束。报告据此给出性能分析、风险评估与后续设计路线，并以最佳 INT8 设计点为基础设计了一个以录制数据回放和逐拍核对为核心的演示方案。

\medskip
\noindent\textbf{关键词：}FPGA；视觉-语言-动作模型；\pizero；脉动阵列；硬核级联；位真验证；设计验证；演示方案
\end{abstract}

\bigskip
\noindent\fbox{\parbox{0.97\textwidth}{\small
\textbf{数据来源与标注约定。}
本报告关于项目的全部事实来自项目简报（\texttt{PI0\_FPGA\_PROJECT\_BRIEF.md}），简报中没有的数字不作为项目结果。
\tM{} 为板上实测或周期精确 RTL 仿真结果；\tP{} 为校准模型预测；\tMsim{} 与 \tMsw{} 为简报原样保留的 RTL 仿真实测与软件模拟实测；
\tD{} 为由简报数字直接算术得到的量，算式随文给出；
\tE{} 为依据模型结构所作的量级估计，不是项目结果；
\tR{} 为本报告提出、尚未实施的设计或验证建议；
\tB{Bxx} 标出\textbf{尚需在 VP815 板卡（或 Jetson AGX Thor 对照平台）上实测}的数据，编号对应第~\ref{sec:board} 节的待上板测量清单，清单中给出每项的测量方法、条件与判据。第~\ref{sec:req}--\ref{sec:impl} 节为设计方案，只描述设计本身，不含此标记；标记只出现在验证、性能、演示与后续计划相关的章节中。
简报未给出的实现细节（例如具体的指令编码、各节点的精确资源数），本报告只描述其设计原则，不给出具体数值。}}

\tableofcontents
\newpage

%======================================================================
\section{引言}\label{sec:intro}
%======================================================================

\subsection{项目目标}

视觉-语言-动作（Vision-Language-Action，VLA）模型把大规模预训练的视觉-语言模型与机器人动作生成结合起来，已成为通用机器人策略的主流形态~\cite{pi0,openvla,rt2}。本项目的目标是把一个完整的 VLA 模型——Physical Intelligence 的 \pizero~\cite{pi0}——的整次推理映射到单片 FPGA 上执行：主机只负责装载输入、启动阵列与读回结果，视觉编码、语言模型前缀、动作专家的全部迭代，以及其间的全部归一化、注意力、激活与量化运算，都在 FPGA 内部完成。

这一目标有三层含义。第一是\emph{完整性}：不把任何一段计算留给主机或其他处理器，从而使时延、数值行为与资源消耗都完全由 FPGA 设计决定，便于分析与优化。第二是\emph{可验证性}：每一次硅片运行的结果都可以与编译器给出的期望内存逐位核对，使设计与调度的每一次改动都有确定的正确性判据。第三是\emph{可度量性}：以同一模型在 NVIDIA Jetson AGX Thor 上的实测时延为外部参照，以相对 fp32 参考实现的机器人关节角误差为精度指标，使结果可以与端侧 GPU 平台直接对比。

\subsection{主要设计挑战}

把 \pizero{} 完整映射到单片 FPGA，需要同时应对以下五个挑战，它们贯穿本报告的各个章节：

\textbf{（1）规模与片上存储的失配。} 模型的 INT8 权重约 2.7~GB \tE，而片上 RAM 约 23~MB \tD，相差两个数量级以上。所有权重都必须从外存按分块流入，外存访问的组织方式直接决定性能。

\textbf{（2）两类形状迥异的负载。} 视觉与前缀 GEMM 的行数为数百，属于算力或激活重流受限；动作专家 GEMM 的行数只有 51，且要重复 10 步，属于权重搬运受限。二者必须共享同一套硬件。

\textbf{（3）非线性与量化运算的比重。} 归一化、softmax、激活、RoPE 与量化相关运算几乎不贡献乘加，却出现在每两个 GEMM 之间。若向量通路没有与矩阵阵列同比例设计，它们会按 Amdahl 定律成为主导项。

\textbf{（4）硬核与布线的约束。} MLP72 的硬级联只在同列相邻单元之间存在，级联连接与乘法模式是静态参数；链之外的连线都要经过通用布线，而布线资源先于乘法器耗尽。

\textbf{（5）正确性与可调试性。} 一个 chunk 包含数万条命令与数十个并发节点，任何一个同步遗漏或数值偏差都可能在最终动作上表现为难以解释的误差。设计必须提供从最终结果回溯到出错指令的手段。

本设计对五个挑战的回应分别是：以 GDDR6 为枢纽、条带化与分块流入（第~\ref{sec:arch}、\ref{sec:compiler} 节）；两种链深与编译器的按深度类分块（第~\ref{sec:chain} 节）；独立向量节点与结果通路 epilogue（第~\ref{sec:vector} 节）；列并行链与 fabric-light 变体（第~\ref{sec:chain}、\ref{sec:impl} 节）；位真模型与逐拍核对（第~\ref{sec:numerics}、\ref{sec:verif} 节）。

\subsection{范围}

本报告覆盖以下内容：需求与约束（第~\ref{sec:req} 节）；负载分析（第~\ref{sec:workload} 节）；总体架构（第~\ref{sec:arch} 节）；GEMM 链与 PV 链设计（第~\ref{sec:chain} 节）；向量节点设计（第~\ref{sec:vector} 节）；数值方案（第~\ref{sec:numerics} 节）；编译器设计（第~\ref{sec:compiler} 节）；FPGA 物理实现（第~\ref{sec:impl} 节）；验证方案与验证结果（第~\ref{sec:verif} 节）；待上板测量清单（第~\ref{sec:board} 节）；性能分析（第~\ref{sec:perf} 节）；演示方案（第~\ref{sec:demo} 节）；风险与后续计划（第~\ref{sec:risk} 节）。关于 VLA 在 FPGA 上部署的文献背景，请参阅配套的综述文档；本报告只在需要时引用相关工作。

\subsection{阅读建议}

本报告按“需求—负载—架构—微结构—数值—编译—实现—验证—性能—演示—风险”的顺序组织。关注硬件设计的读者可重点阅读第~\ref{sec:arch}、\ref{sec:chain}、\ref{sec:vector} 与 \ref{sec:impl} 节；关注验证的读者可重点阅读第~\ref{sec:numerics}、\ref{sec:compiler} 与 \ref{sec:verif} 节；只需了解结论与演示的读者可阅读摘要、第~\ref{sec:perf}、\ref{sec:demo} 与 \ref{sec:risk} 节。附录~\ref{app:checklist} 汇总了评审检查项。全文所有项目数字都带有来源标注，阅读时请注意区分实测（\tM）、模型预测（\tP）与本报告的推算（\tD、\tE）和建议（\tR）。

\subsection{术语}

\begin{itemize}
  \item \textbf{chunk}：\pizero{} 的一次完整推理，输出一个含 50 个未来动作的动作块。
  \item \textbf{前缀（prefix）}：图像 token 与语言 token 组成的序列，经 Gemma-2B 计算一次，产生 KV cache。
  \item \textbf{动作专家（action expert）}：约 3 亿参数的 Gemma 结构网络，对 51 个 token 迭代 10 个 Euler 积分步。
  \item \textbf{MLP72}：Speedster7t 的机器学习处理器硬核块；\textbf{BRAM72K} 与 \textbf{LRAM2K} 为片上 RAM 硬核。
  \item \textbf{NAP}：fabric 逻辑进入二维片上网络（NoC）的接入点。
  \item \textbf{链（chain）}：沿 MLP72 硬级联串接的一组 MLP72，构成一个一维矩阵乘阵列；\textbf{feeder} 是链底部读入激活的 MLP72。
  \item \textbf{epilogue}：在链的结果通路上、结果离开节点之前完成的后处理（如反量化与激活）。
  \item \textbf{位真（bit-exact）}：硅片结果与编译器期望内存逐位一致。
\end{itemize}

%======================================================================
\section{需求与约束}\label{sec:req}
%======================================================================

\subsection{功能需求}

\textbf{F1　完整推理。} 输入为两路相机图像、语言指令 token 与机器人状态；输出为 50 个未来动作组成的动作块。中间过程包括：SigLIP So400m/14 视觉编码（两路图像各 256 个 token）；由图像与语言 token 组成的 525 个 token 前缀经 Gemma-2B 计算，得到 KV cache；动作专家对 51 个 token（1 个状态 + 50 个动作）迭代 10 个 Euler 积分步，每步关注 525 个前缀键与自身的键。

\textbf{F2　主机最小参与。} 主机只装载输入、启动阵列、读回结果；FPGA 内部的计算、调度与同步不依赖主机干预。

\textbf{F3　确定性结果。} 同一输入、同一程序在硅片上的每次运行，其结果必须与编译器期望内存逐位一致。

\textbf{F4　多设计点支持。} 同一套编译器应能为不同的硬件设计点（INT8 链、INT4 链、带 epilogue 的链、fabric-light 链等）生成程序，以支持设计空间探索。

\subsection{输入输出接口}

主机与阵列之间的数据接口包括 \tR：
\begin{itemize}
  \item \textbf{图像输入}：两路相机图像。SigLIP So400m/14 的图块大小为 14，每路 256 个 token 对应 $16\times16$ 个图块，即每路约 $224\times224$ 像素 \tE；图像预处理（缩放、归一化）在主机侧完成还是在阵列内完成，应在接口定义中明确。
  \item \textbf{语言输入}：已分词的指令 token 序列；与图像 token 合计构成 525 个 token 的前缀。
  \item \textbf{状态输入}：机器人当前状态，作为动作专家的 1 个状态 token 的来源。
  \item \textbf{动作输出}：50 个未来动作组成的动作块。
  \item \textbf{控制与状态}：启动标志、完成标志，以及可选的错误与调试信息区域。
\end{itemize}
接口定义应独立于设计点，由编译器以固定的地址约定输出，使主机程序在更换比特流后不需修改。

\subsection{精度需求}

检查点为在 DROID~\cite{droid} 数据集上微调的 \pizero；精度评估在 46 帧留出的录制数据上进行，以 fp32 PyTorch 参考实现为基准，报告 7 自由度机械臂在整个动作块上关节角的最大误差与平均误差。精度需求以“可用”为界：最佳 INT8 设计点的关节最大误差 $0.665^\circ$、平均 $0.072^\circ$ \tM{} 被视为可用；INT4 设计点的关节最大误差 $44.9^\circ$ \tM{} 被判定为不可用。本报告建议把这一“可用”界线显式化为验收阈值，并补充闭环指标（第~\ref{sec:verif} 节）\tR。

\subsection{性能需求}

外部参照为 NVIDIA Jetson AGX Thor 上以 bf16 编译的同一 \pizero{} chunk：125.3~ms \tM。项目以“在匹配精度下逼近该参照”为长期目标；当前最佳可用设计点为 1.221~s \tM，约为参照的 9.7 倍 \tD（$1.221/0.1253\approx9.74$）。需要说明，简报第 6 节写作“约 8 倍”，按简报自身数字，8 倍对应的是精度不可用的 INT4 设计点（$1.007/0.1253\approx8.04$ \tD）；本报告按匹配精度口径采用 9.7 倍。

\subsection{验收标准}

为使每一个新设计点的评估有统一的判据，本报告建议采用以下验收标准 \tR：
\begin{enumerate}
  \item \textbf{正确性}：在完整 35 节点阵列上运行整 chunk，全部 46 帧的外存结果与编译器期望内存逐位一致；
  \item \textbf{精度}：46 帧上的关节最大误差与平均误差不劣于当前最佳 INT8 设计点（$0.665^\circ$ 与 $0.072^\circ$ \tM），或在给出任务层面的理由后以新的阈值替代；
  \item \textbf{性能}：整 chunk 时延优于当前最佳可用设计点（1.221~s \tM），并给出与校准模型预测的偏差；
  \item \textbf{可复现}：构建产物完整可追溯，任何人可以用同一组产物复现上述结果；
  \item \textbf{稳定性}：连续多次运行时延的波动在可接受范围内，且无挂死与偶发不一致。
\end{enumerate}
第 5 条目前缺少数据，需要以演示方案中的批量回放（第~\ref{sec:demo} 节）或专门的稳定性测试补齐。

\subsection{非功能需求}

除上述功能、精度与性能需求外，设计还需满足以下非功能需求：

\textbf{N1　可调试性。} 当硅片结果出错时，必须能够定位到出错的节点与指令。本设计以“逐拍核对”满足这一需求：编译器给出每一次外存写入的期望值，任何不一致都可以回溯到其生产者。

\textbf{N2　可演进性。} 模型结构、分块与调度的变化不应要求重新综合；只有当硬件结构本身改变（例如新增 epilogue 或改变链深）时才需要新的比特流。本设计以“程序即模型”的原则满足这一需求。

\textbf{N3　可度量性。} 每个设计点都应有可复现的时延、资源与精度数据，并有一个能解释这些数据的性能模型。本设计以校准模型与周期精确仿真满足这一需求。

\textbf{N4　可移植的软件接口。} 主机侧接口应与具体设计点无关：无论比特流采用何种链或时钟，主机都以相同的方式装载输入、启动与读回。这使演示、评测与后续集成不受硬件迭代的影响。

\subsection{需求追踪}

表~\ref{tab:trace} 把需求映射到设计要素与验证方法，作为后续章节的索引。

\begin{table}[h]
\centering
\caption{需求—设计—验证追踪。}
\label{tab:trace}
\small
\begin{tabularx}{\textwidth}{l Y Y}
\toprule
需求 & 设计要素 & 验证方法 \\
\midrule
F1 完整推理 & 节点阵列 + 向量节点覆盖全部算子（第~\ref{sec:arch}--\ref{sec:vector} 节） & 整 chunk 位真；46 帧精度 \\
F2 主机最小参与 & 自主节点、GDDR6 标志同步、主机接口（第~\ref{sec:arch} 节） & 主机开销测量；接口测试 \\
F3 确定性结果 & 静态调度、位真模型、期望内存（第~\ref{sec:numerics}、\ref{sec:compiler} 节） & L1--L4 逐位核对 \\
F4 多设计点 & 硬件描述驱动的编译器（第~\ref{sec:compiler} 节） & 各设计点整 chunk 位真 \\
精度 & W8A8 + SmoothQuant + uint8 概率 + bf16 接口 & L0 对 L1 的动作误差 \\
性能 & 列并行链、epilogue、调度优化 & 实测时延；校准模型对照 \\
N1 可调试性 & 期望内存 + 故障定位 & 注入错误的定位测试 \\
N2 可演进性 & 程序即模型 & 同一比特流上更换调度 \\
N3 可度量性 & 校准模型 + 周期精确仿真 & 模型预测与实测偏差 \\
N4 接口稳定 & 设计点无关的主机协议 & 多比特流下的同一主机程序 \\
\bottomrule
\end{tabularx}
\end{table}

\subsection{平台约束}

表~\ref{tab:platform} 汇总目标平台的资源与约束。

\begin{table}[h]
\centering
\caption{目标平台（依据简报）。}
\label{tab:platform}
\small
\begin{tabularx}{\textwidth}{l Y}
\toprule
项目 & 内容 \\
\midrule
器件 & Achronix Speedster7t AC7t1500 \\
板卡 & VP815 PCIe 卡，PCIe Gen4 x16 \\
外存 & 16 路 $\times$ 2~GB GDDR6，经二维 NoC 访问 \\
片上网络 & 二维 NoC；fabric 逻辑经 NAP 接入 \\
矩阵硬核 & 2,560 个 MLP72：每周期 16 个 INT8$\times$INT8 乘加，或 SIGNED 4$\times$4 模式下 32 个有符号 INT4$\times$INT4 乘加；块浮点模式；累加器；把一个操作数字转发给同列下一个 MLP72 的硬级联 \\
片上 RAM & 2,560 块 BRAM72K；LRAM2K 小 RAM（每个 MLP72 占用与之配对的 LRAM2K） \\
静态参数 & 级联连接与乘法器模式是比特流的静态参数 \\
工具链 & Achronix ACE 10.5.2 + Synplify；Verilator 用于 RTL 仿真 \\
\bottomrule
\end{tabularx}
\end{table}

其中有两条约束对架构影响最大。其一，\emph{级联连接与乘法器模式是静态的}：一条链的长度、每个 MLP72 的乘法模式（INT8、INT4 或块浮点）在综合时确定，运行时只能在预先布设的备选路径之间切换。其二，\emph{片上存储远小于模型}：BRAM72K 若按名义容量每块 72~Kb 计，合计约 23~MB \tD（$2560\times72~\text{Kb}\div8$），而 \pizero{} 的 INT8 线性层权重约 2.7~GB \tE（见第~\ref{sec:workload} 节），权重必须存放在 GDDR6 中并按分块流入。

%======================================================================
\section{负载分析}\label{sec:workload}
%======================================================================

\subsection{模型结构}

按简报，\pizero{} 由三部分组成：
\begin{itemize}
  \item \textbf{视觉编码器}：SigLIP So400m/14~\cite{siglip}，27 层，宽 1152，16 头 $\times$ 72；
  \item \textbf{语言模型}：Gemma-2B~\cite{gemma}，18 层，宽 2048，MLP 宽 16384，采用 GeGLU，8 个查询头 $\times$ 256，共享 1 个 KV 头；
  \item \textbf{动作专家}：约 3 亿参数，Gemma 结构，18 层，宽 1024，MLP 4096，关注 VLM 的 KV cache。
\end{itemize}
VLM 部分即 PaliGemma~\cite{paligemma}。动作由流匹配生成：每次推理 10 个 Euler 积分步。

\subsection{GEMM 形状}

表~\ref{tab:gemm} 给出按上述结构参数推得的每层主要 GEMM 形状（$M$ 为行数即 token 数，$K$ 为输入维，$N$ 为输出维）。其中两处参数简报未给出，按公开结构推断并标为 \tE：SigLIP So400m 的 MLP 宽度取 4304；动作专家的注意力需与前缀 KV（1 个 KV 头、头宽 256）兼容，其查询投影取 8 头 $\times$ 256。

\begin{table}[t]
\centering
\caption{每层主要 GEMM 形状（依据简报结构参数推得；标 \tE{} 的维度为推断）。}
\label{tab:gemm}
\small
\begin{tabular}{l l r r r r}
\toprule
阶段 & GEMM & $M$ & $K$ & $N$ & 每 chunk 执行次数 \\
\midrule
\multirow{3}{*}{SigLIP（27 层）} & QKV 投影 & 512 & 1152 & 3456 & 27 \\
 & 输出投影 & 512 & 1152 & 1152 & 27 \\
 & MLP 上/下投影 & 512 & 1152 / 4304\tE & 4304\tE / 1152 & 27 \\
\midrule
\multirow{5}{*}{Gemma-2B 前缀（18 层）} & Q 投影 & 525 & 2048 & 2048 & 18 \\
 & K、V 投影 & 525 & 2048 & 256 + 256 & 18 \\
 & 输出投影 & 525 & 2048 & 2048 & 18 \\
 & gate、up 投影 & 525 & 2048 & 16384 $\times$ 2 & 18 \\
 & down 投影 & 525 & 16384 & 2048 & 18 \\
\midrule
\multirow{5}{*}{动作专家（18 层 $\times$ 10 步）} & Q 投影 & 51 & 1024 & 2048\tE & 180 \\
 & K、V 投影 & 51 & 1024 & 256 + 256 & 180 \\
 & 输出投影 & 51 & 2048\tE & 1024 & 180 \\
 & gate、up 投影 & 51 & 1024 & 4096 $\times$ 2 & 180 \\
 & down 投影 & 51 & 4096 & 1024 & 180 \\
\bottomrule
\end{tabular}
\end{table}

按表~\ref{tab:gemm} 累加，每个 chunk 的线性层 MAC 约为：Gemma-2B 前缀 1.04~TMAC，SigLIP 0.21~TMAC，动作专家 0.16~TMAC，另有数十 GMAC 的注意力矩阵乘，合计约 1.4 至 1.5~TMAC \tE，与简报的 1.47~TMAC \tP{} 一致。两类形状的对比是整个设计的出发点：简报指出，前缀 GEMM 的 $M\approx525$，属于算力受限，激活的反复重流占主导；动作专家 GEMM 的 $M=51$，属于权重搬运受限。

\subsection{注意力}

前缀注意力：8 个查询头共享 1 个 KV 头，头宽 256，序列长 525，每层的 $QK^\top$ 与 $PV$ 各约 $525\times525\times256\times8$ 次乘加。动作专家注意力：51 个查询 token 关注 525 个前缀键与 51 个自身键，共 576 个键。按简报的剪除规则，前缀最后一层只需计算其 K/V，因为其后续输出不被动作专家使用。注意力概率以 uint8 表示，“概率 $\times$ 值”（$PV$）由专门的 PV 链计算（uint8 $\times$ int8）。

\subsection{非线性与量化相关运算}

按简报，向量节点负责以下运算：LayerNorm/RMSNorm、softmax、GeGLU、GELU、RoPE、残差相加、按 token 的量化（行尺度在运行时计算）、按“行尺度 $\times$ 列尺度”的反量化，以及融合的量化尾。超越函数（GELU、exp、rsqrt、sigmoid）基于查表实现。这些运算几乎不贡献乘加，但它们位于每两个 GEMM 之间的关键路径上，其总量随层数与积分步数线性增长。按简报的校准模型，它们在最佳 INT8 设计点上约占 50\% 的时延 \tP。

\subsection{存储需求}

\begin{itemize}
  \item \textbf{权重}：INT8 下，VLM（SigLIP + Gemma-2B）线性层权重约 2.4~GB，动作专家约 0.31~GB \tE；二者之和远超片上容量，全部存放在 GDDR6。动作专家权重在 10 个积分步中各被完整读取一次，若不能驻留片上，每个 chunk 的专家权重流量约 3.1~GB \tE，超过前缀权重流量。
  \item \textbf{KV cache}：前缀 KV 规模为 $525\times18\times2\times256$ 个元素，bf16 下约 9.7~MB、INT8 下约 4.8~MB \tE，每个积分步的每一层都要重读。
  \item \textbf{激活}：前缀阶段最大的中间激活为 gate/up 投影输出（$525\times16384$），INT8 下约 8.6~MB \tE；动作专家阶段的激活规模小一个数量级以上。
\end{itemize}
这些量决定了一个基本设计选择：\emph{权重、激活与 KV 都以 GDDR6 为主存储，片上 BRAM 只作为分块缓冲}。节点之间的中间结果也经 GDDR6 交换（第~\ref{sec:arch} 节）。

\subsection{按阶段的负载特征}

\textbf{视觉编码阶段（SigLIP）。} 两路图像共 512 个 token，27 层，宽 1152。其 GEMM 的 $M=512$，与前缀相近，属于算力或激活重流受限；注意力为 16 头 $\times$ 72，头宽较小，$QK^\top$ 的 $K$ 维只有 72，在列并行链上属于“小 $K$”运算（第~\ref{sec:chain} 节的行周期下限问题）。两路图像的编码彼此独立，可以在节点之间并行。

\textbf{前缀阶段（Gemma-2B）。} 525 个 token，18 层。线性层权重约 1.98 G 个参数 \tE，是整个 chunk 中权重量最大、MAC 最多的阶段（约 1.04~TMAC \tE）。MLP 的 gate 与 up 投影输出宽 16384，GeGLU 需要把两者逐元素相乘后再进入 down 投影；down 投影的 $K=16384$ 是全 chunk 最长的归约维，按简报分为 4 段计算。前缀的注意力采用 1 个 KV 头，K 与 V 投影的输出只有 256 维，这使 KV cache 很小，但也使 K/V 投影的 $N$ 很窄，在 32 级链上需要注意边角浪费。

\textbf{动作专家阶段。} 51 个 token，18 层，重复 10 步。每步的线性层 MAC 约 16~GMAC \tE，但权重约 0.31~GB（INT8）\tE 需要在每步重新流过链。每步的每一层还要读取该层的前缀 KV（525 个键）并与自身 51 个键拼接。步与步之间有严格的数据依赖，只能流水、不能并行。

\subsection{数据搬运预算}

以 INT8 计，每个 chunk 的主要外存流量可以分解为 \tE：
\begin{itemize}
  \item 权重：前缀与视觉约 2.4~GB（至少一遍），动作专家约 3.1~GB（10 遍）；
  \item 激活重流：在列并行链上，每个 $N$ 方向的分块都要把全部 $M$ 行激活重读一遍；对前缀的宽 GEMM（如 $N=16384$），重读次数与分块数成正比，是前缀阶段的主要外存流量来源之一；
  \item 中间结果往返：链的 int32 结果与向量节点之间的 bf16 数据经 GDDR6 往返；epilogue 的作用之一就是削减这部分流量；
  \item KV 读取：前缀 KV 每步每层读取一次，10 步共约 48~MB（INT8）或 97~MB（bf16）\tE。
\end{itemize}
这一预算说明，“把哪些数据留在片上”是比“用多少乘法器”更根本的设计问题。在片上容量约 23~MB \tD{} 的约束下，前缀 KV 是最值得驻留的候选，其次是动作专家的部分权重（第~\ref{sec:risk} 节）。

\subsection{各类 GEMM 的算术强度}

若权重每次从外存读取，一个 $M\times K\times N$ 的 GEMM 在 INT8 下的算术强度上限约为 $2M$ ops/byte（$M\ll K,N$ 时）。据此，前缀与视觉 GEMM 约 1,000~ops/byte，动作专家 GEMM 约 100~ops/byte \tE。这一差别在硬件上表现为：前缀 GEMM 只要片上缓冲能容纳足够大的权重分块，就能让链保持计算受限；动作专家 GEMM 则无论链多快，都受制于每步重新装载权重的速度。两类负载因此需要不同的优化手段，而它们必须共享同一个节点阵列。

\subsection{负载对硬件的需求小结}

综合本节的分析，\pizero{} 的负载对硬件提出以下需求，它们直接对应第~\ref{sec:arch} 至第~\ref{sec:vector} 节的设计选择：
\begin{enumerate}
  \item \textbf{高效的 INT8 矩阵乘}：约 1.47~TMAC/chunk \tP{} 的乘加必须由硬核承担，且在 $M$ 为数百与 $M=51$ 两类形状上都保持可接受的效率；
  \item \textbf{与外存之间的高效分块流动}：权重、激活与 KV 都以外存为主存储，分块的装入必须能与计算重叠；
  \item \textbf{足够的向量吞吐}：每两个 GEMM 之间都有归一化、激活或量化，向量吞吐必须与矩阵吞吐匹配；
  \item \textbf{注意力的专门支持}：两个操作数都为运行时数据、且概率为无符号格式的矩阵乘；
  \item \textbf{精确可复现的数值行为}：整数运算、查表与舍入必须完全确定，以支持逐位核对。
\end{enumerate}

\subsection{向其他 VLA 模型的适配}

由于本架构以“程序即模型”为原则，适配其他结构相近的 VLA 模型主要是编译器的工作。以几类模型为例 \tR：
\begin{itemize}
  \item \textbf{同属流匹配、结构相近的模型}（如 \pizero{} 的后续版本~\cite{pi05}）：若骨干仍为 PaliGemma 类结构、动作专家仍为 Gemma 类结构，主要变化在于层数、宽度、token 数与参数，编译器可以直接处理；若引入新的归一化形式（例如带条件输入的归一化），则需要在向量节点中新增运算，并同步扩展位真模型。
  \item \textbf{小型 VLA}（如 SmolVLA~\cite{smolvla}）：参数量小一个数量级，权重可能有更大比例驻留片上，动作专家的权重装载瓶颈会显著缓解；这类模型可能更适合在当前器件上达到更低的时延。
  \item \textbf{扩散动作头的模型}（如 GR00T N1~\cite{gr00t} 的扩散 Transformer）：推理结构与流匹配同构，都是对同一网络的多次迭代，区别在于步数与更新公式，更新公式由向量节点实现。
  \item \textbf{离散 token 自回归的模型}（如 OpenVLA~\cite{openvla}）：解码阶段 $M=1$，属于极端的带宽受限负载，本架构的列并行链在 $M=1$ 时权重复用度为 1，效率会很低；这类模型不是本架构的目标。
\end{itemize}

%======================================================================
\section{总体架构}\label{sec:arch}
%======================================================================

\subsection{设计原则}

总体架构遵循四条原则：
\begin{enumerate}
  \item \textbf{计算单元自主}：每个计算单元（节点）从外存取自己的程序并自行推进，不依赖中央调度器；依赖关系由编译器在程序中显式表达。
  \item \textbf{矩阵与向量分离}：整数乘加在 MLP72 级联链上完成，浮点的逐元素与归约运算在独立的向量节点上完成，二者以 bf16 接口经外存交换数据；在收益明确处，再把部分向量运算并入链的结果通路（epilogue）。
  \item \textbf{外存为枢纽}：所有节点通过各自的 NAP 访问 GDDR6，中间结果与同步标志都放在 GDDR6 中，避免在 fabric 中布设大型交叉开关。
  \item \textbf{程序即模型}：模型结构、分块、调度与数值参数都体现在编译器生成的程序中；更换模型或调度只需重新编译，不需重新综合。
\end{enumerate}

\subsection{节点阵列}

\begin{figure}[t]
\centering
\begin{tikzpicture}[font=\footnotesize, >=Stealth,
  nd/.style={draw, rounded corners=1.5pt, minimum height=8mm, align=center, fill=#1, text width=2.4cm},
  nd/.default=white]
\node[draw, fill=gray!10, minimum width=13cm, minimum height=7mm] (noc) {二维片上网络（NoC）：每个节点经独立 NAP 接入};
\node[draw, fill=yellow!15, minimum width=13cm, minimum height=7mm, below=18mm of noc] (gddr) {GDDR6（16 $\times$ 2 GB）：程序 \quad 权重（条带化） \quad 激活 \quad KV \quad 标志字（WAIT/POST）};
\node[nd=blue!10] (a) at ($(noc.south west)+(1.7cm,-9mm)$) {GEMM 链 $\times$12\\32 级乘法};
\node[nd=blue!10, right=0.8cm of a] (b) {GEMM 链 $\times$12\\16 级乘法};
\node[nd=green!10, right=0.8cm of b] (c) {PV 链 $\times$3\\16 级，uint8$\times$int8};
\node[nd=orange!15, right=0.8cm of c] (d) {向量节点 $\times$8\\各 2 lane，bf16};
\foreach \n in {a,b,c,d} { \draw[<->] (\n.north) -- (\n.north |- noc.south); \draw[<->] (\n.south) -- (\n.south |- gddr.north); }
\node[draw, fill=gray!5, right=4mm of noc, text width=1.5cm, align=center] (host) {主机\\PCIe Gen4};
\draw[<->] (host) -- (noc);
\end{tikzpicture}
\caption{节点阵列的组织：35 个自主节点经各自的 NAP 访问 GDDR6；节点之间以 GDDR6 中的标志字同步，主机只装载输入、启动与读回。}
\label{fig:array}
\end{figure}

阵列由 35 个自主节点组成（图~\ref{fig:array}）：
\begin{itemize}
  \item 24 个整数 GEMM 链节点：12 个 32 级乘法链、12 个 16 级乘法链；
  \item 3 个 PV 链节点：负责注意力概率 $\times$ 值，uint8 $\times$ int8，16 级；
  \item 8 个向量节点：各 2 条 lane。
\end{itemize}
每个节点拥有自己的 NAP，从 GDDR6 取自己的程序。

两种链深的并存是针对 GEMM 形状多样性的静态回应：32 级链在 $N$ 方向一次产出更多输出列，适合大 $N$ 的 GEMM；16 级链的行周期下限更低，适合 $K$ 较短或需要更细粒度负载均衡的 GEMM。编译器按“深度类”把每个 GEMM 分配给合适的链（第~\ref{sec:compiler} 节）。

\subsection{备选架构与取舍}

在确定当前架构之前，需要在几种典型的 FPGA 加速器组织之间做出取舍。本节从 VLA 负载的特点出发比较四种方案。

\textbf{方案 A：单一大阵列 + 中央控制器。} 把尽可能多的 MLP72 组织成一个二维脉动阵列，由一个中央控制器按层调度，非线性运算由一个与阵列相连的向量单元完成。其优点是控制简单、激活可以在阵列内广播复用；缺点有三：一是 MLP72 的硬级联只在同列相邻单元之间存在，二维阵列的另一个方向必须经通用布线，布线压力随阵列规模超线性增长；二是 $M=51$ 的动作专家 GEMM 在大阵列上利用率很低，而阵列无法被拆开服务不同的 GEMM；三是中央控制器必须覆盖全部阵列，其控制通路本身成为布线与时序的瓶颈。

\textbf{方案 B：逐层空间流水（模型专用数据流）。} 为每一层或每一类算子实例化专门的硬件单元，数据在单元之间直接流动，尽量不经外存。这类设计在小模型上能取得最高效率，但 \pizero{} 的权重与激活远超片上容量，逐层实例化意味着同一层的硬件在 chunk 的大部分时间里闲置；而且模型每次变化都需要重新综合，违背需求 N2。

\textbf{方案 C：软核处理器 + 硬核阵列（overlay）。} 在 FPGA 上实现一个或数个通用软核处理器，控制硬核矩阵阵列与向量单元。其优点是可编程性最好，缺点是软核的指令获取与控制开销会进入关键路径，而且少数几个软核很难同时驱动数十条独立的链。

\textbf{方案 D（本设计）：多节点自主阵列。} 把阵列拆成数十个小而独立的节点，每个节点有自己的程序、自己的 NoC 接入点，节点之间只通过外存交换数据与同步。它在上述三者之间取了折中：
\begin{itemize}
  \item 与方案 A 相比，每条链只沿硬级联一维延伸，节点之间不需要 fabric 中的大型互连；小节点可以各自服务不同形状的 GEMM，由编译器在节点之间分配分块，相当于在软件中完成阵列的“拆分与合并”。
  \item 与方案 B 相比，硬件单元不绑定到具体层，而是在整个 chunk 中被反复复用；模型变化只需重新编译。
  \item 与方案 C 相比，每个节点的控制逻辑很小，只需按顺序执行自己的命令流，不需要通用处理器。
\end{itemize}
其代价是：节点间所有数据与同步都经外存，外存访问时延会进入关键路径；每个节点都需要自己的馈送与引出通路，仍然消耗布线资源。当前实测中“feeder 停顿约 6\%”\tP{} 与“布线是约束资源”都是这一代价的体现。

\subsection{节点的功能组成}

简报没有给出节点内部的具体实现，本节只从功能上说明一个节点需要具备的组成部分，作为设计与评审的检查清单：
\begin{itemize}
  \item \textbf{程序获取}：经本节点的 NAP 从 GDDR6 读取命令流，按序执行；
  \item \textbf{外存访问}：发起权重、激活、KV 与标志字的读写请求，处理 NoC 返回的数据；
  \item \textbf{本地缓冲}：GEMM 链节点需要激活缓冲（feeder 侧，双缓冲）与各级权重缓冲（BRAM72K）；向量节点需要行缓冲以支持归约运算；
  \item \textbf{同步}：执行 \texttt{WAIT}（轮询标志字）与 \texttt{POST}（写标志字）；
  \item \textbf{结果引出}：GEMM 链节点把各级累加器的结果经 epilogue（若有）写回外存；
  \item \textbf{跨时钟域}：节点的计算部分在阵列域或向量域运行，程序获取与 NoC 接口在 fabric 域运行，二者之间以异步 FIFO 连接。
\end{itemize}

\subsection{命令流的抽象}

简报只明确给出了两条同步命令（\texttt{WAIT} 与 \texttt{POST}）。从编译器与节点的接口看，一个节点的命令流至少需要表达以下几类操作：装载权重分块到链的各级 BRAM；按行流过激活并产出结果；选择链的运行模式（例如融合链对、链分裂、epilogue 类型）；执行一段向量运算序列；以及同步。本报告不给出具体的指令编码。需要强调的设计原则是：命令流中不存在数据依赖的分支，所有循环次数与地址在编译期确定。这正是期望内存可以被静态计算、硅片结果可以被逐拍核对的前提。

\subsection{外存组织}

GDDR6 在本设计中承担四类数据：程序、权重、运行时数据（激活、KV、中间结果）与同步标志。四类数据的访问模式不同：程序按节点顺序读取；权重按分块顺序读取，每个 chunk 读一遍（专家权重读 10 遍）；运行时数据读写交替；标志字被高频轮询。编译器的条带化地址映射把大块的权重与激活分散到多个通道上，使多个节点的并发访问不集中于同一通道；标志字的放置则应避免与大流量数据竞争同一通道，以降低同步时延 \tR。

\subsection{设计不变量}

为保证逐拍核对始终有效，设计需要维持以下不变量，任何硬件或编译器的改动都不得破坏它们：
\begin{enumerate}
  \item \textbf{无数据依赖的控制流}：命令流中所有循环次数、地址与模式选择都在编译期确定；
  \item \textbf{显式同步}：所有跨节点的数据依赖都由 \texttt{WAIT}/\texttt{POST} 保护，节点之间不存在隐式的时序假设；
  \item \textbf{确定性数值}：所有运算（包括查表与舍入）对同一输入产生同一输出，不依赖执行顺序或时序；
  \item \textbf{单一数值定义}：硬件、位真模型与期望内存使用同一份数值定义。
\end{enumerate}
这四条不变量是“确定性结果”（F3）与“可调试性”（N1）两项需求在设计层面的具体化。运行时可重构（融合链对、链分裂）也必须在编译期选定模式，以维持第 1 条。

\subsection{同步机制}

节点之间以 GDDR6 中的标志字同步，程序中有两条同步命令：\texttt{POST} 在某个数据生产完成后写入标志字，\texttt{WAIT} 在消费数据前轮询标志字直到其满足条件。编译器在展开数据依赖时为每一对生产者—消费者插入 \texttt{POST}/\texttt{WAIT}。这一机制的优点是：没有中央控制器；同步语义简单、可在编译器的期望内存模型中精确建模；节点数可以扩展而不需要扩展控制通路。其代价是每次同步都要经过 GDDR6 往返，这部分开销体现在简报所称的“feeder 停顿”中（约 6\% \tP）。

为避免同步开销放大，编译器采用两项措施：一是把权重装载提升到 \texttt{WAIT} 之前，使等待期间可以预取不依赖上游结果的数据；二是以注意力行块与层间重叠等手段，让一个节点在等待某个依赖时有其他工作可做。

\subsection{一个 chunk 的执行过程}

为把上述各部分串联起来，本节按时间顺序描述一个 chunk 在节点阵列上的执行过程。这一描述是依据简报给出的架构与编译优化所作的概括，具体的节点分配与顺序由编译器决定。

\textbf{第一步：输入就绪。} 主机把两路图像、语言 token 与机器人状态写入 GDDR6 的约定区域，然后写入启动标志。所有节点从 GDDR6 读取各自的程序，第一条相关命令通常是等待启动标志。

\textbf{第二步：视觉编码。} 两路图像的图块嵌入进入 SigLIP 的 27 层。每层中，GEMM 链节点承担 QKV 投影、输出投影与 MLP 的线性层，PV 链节点承担注意力的“概率 $\times$ 值”，向量节点承担 LayerNorm、softmax、GELU、残差与量化/反量化。两路图像彼此独立，编译器可以把它们的工作交错分配，以提高节点利用率。每一层的输出经 GDDR6 传给下一层，层与层之间以 \texttt{WAIT}/\texttt{POST} 同步；注意力行块与层间重叠使相邻层的部分工作可以并行。

\textbf{第三步：前缀计算。} 视觉 token 与语言 token 拼成 525 个 token 的前缀，进入 Gemma-2B 的 18 层。每层的 K、V 投影结果写入 GDDR6 中的 KV 区域，供后续所有积分步使用。MLP 的 down 投影按 4 段计算并合并部分和。最后一层只计算 K/V，其余部分被编译器剪除，因为动作专家只需要前缀的 KV。

\textbf{第四步：动作专家迭代。} 动作专家对 51 个 token 执行 10 个积分步。每一步的每一层中：GEMM 链节点从 GDDR6 流入该层的专家权重并完成线性层；PV 链节点读取该层的前缀 KV 与本步自身的 KV，完成注意力；向量节点完成归一化、softmax、GeGLU、残差与量化。一步结束时，向量节点完成 Euler 更新，得到下一步的输入。由于步间依赖，第 $k+1$ 步必须等待第 $k$ 步的 Euler 更新完成；编译器能做的是在一步之内的不同层之间重叠，以及提前装载下一层的权重。

\textbf{第五步：输出。} 第 10 步结束后，动作块写入 GDDR6 的输出区域，最后一个节点写入完成标志。主机检测到完成标志后读回动作块；在验证模式下，主机还读回全部中间结果区域，与期望内存逐位比较。

从这一过程可以看出，整个 chunk 中没有任何一次主机干预，也没有任何数据依赖的分支；每个节点在什么时刻做什么，完全由编译期生成的命令流与同步决定。

\subsection{主机接口}

主机经 PCIe Gen4 x16 与板卡通信，职责仅限于：把输入（图像、语言 token、状态）写入 GDDR6 的约定区域；写入启动标志；等待完成标志；读回动作块。程序与权重在初始化阶段一次性装入 GDDR6，之后各次推理复用。简报给出的主机开销约 47~ms/次调用；在当前 1.221~s 的时延下占比约 3.8\% \tD（$47/1221$），但在以百毫秒为目标的后续阶段中将变得显著，需要单独优化（第~\ref{sec:risk} 节）。

\subsection{可扩展性}

本架构在三个方向上具有可扩展性，每个方向都有明确的限制。

\textbf{节点数。} 由于没有中央控制器，增加节点只需在 NoC 上增加接入点、在编译器的硬件描述中增加条目。限制来自三方面：布线（每个节点都有自己的馈送、引出与控制逻辑）、NAP 的数量与分布，以及外存带宽（更多节点意味着更多并发访问）。当前实测表明，布线是最先到达的限制。

\textbf{节点类型的配比。} 链节点与向量节点的数量比可以在综合时调整。当前的 24:3:8（GEMM 链：PV 链：向量节点）配比偏向矩阵侧；第~\ref{sec:risk} 节建议的再配比就是沿这一方向扩展。

\textbf{多器件。} 由于节点间只经外存通信，理论上可以把一部分节点放到另一块 FPGA 上，只要两块板卡之间有共享的存储或足够快的数据通道。一个自然的划分是把视觉与前缀放在一块器件、把动作专家放在另一块器件，并让后者常驻专家权重。这一方向目前没有实测数据，属于远期选项 \tR。

\subsection{外存布局示例}

以一个线性层为例说明外存布局的设计考虑 \tR。权重按“分块优先”的顺序存放：同一分块内各级链需要的列权重连续存放，使一次分块装载成为若干段连续的突发读；不同分块按编译器的执行顺序排列，使相邻的装载请求落在不同的 GDDR6 通道上。激活按行存放，每行的 INT8 数据后附其行尺度，使 feeder 一次读取即可得到一行的数据与尺度。输出结果按 epilogue 的输出格式存放：若 epilogue 已完成反量化与激活，则直接以 bf16 存放，供下游向量节点读取。KV 区域按层与头组织，使动作专家每一步读取某层 KV 时是一段连续访问。这种布局的共同原则是：让每个节点的访问尽量是长的顺序突发，并让并发的节点访问分散在不同通道上。

\subsection{失效模式与容错}

在机器人部署场景中，加速器的失效模式需要被明确考虑 \tR。本设计可能出现的失效包括：某节点因同步错误而永远等待（挂死）；外存读写出错导致结果错误；以及时序裕量不足导致偶发的计算错误。对应的防护措施可以是：主机侧为每次推理设置超时，超时后复位阵列并报告；在输出区域附加校验信息（例如对动作块的简单校验和），使主机能够发现明显的输出损坏；以及在部署前以长时间的压力测试确认时序裕量。由于本设计的每一次推理都有确定的期望行为，任何偏离都可以被快速识别，这比非确定性的执行平台更有利于建立可靠的故障检测。

\subsection{时钟域}

设计分为三个时钟域：\emph{阵列域}驱动 MLP72 链；\emph{向量域}驱动向量节点；\emph{fabric 域}驱动控制逻辑与 NoC 侧接口。简报以“阵列 / 向量 / fabric MHz”记录各设计点的时钟，最佳 INT8 设计点为 500 / 288.9 / 216.7~MHz。三域分离的理由是：MLP72 硬核与级联通路可以运行在远高于通用逻辑的频率；向量节点含浮点与查表逻辑，时序更紧；控制与 NoC 接口逻辑分布广、布线长，频率最低。跨域数据经异步 FIFO 传递，跨域的验证要求见第~\ref{sec:verif} 节。

%======================================================================
\section{GEMM 链与 PV 链设计}\label{sec:chain}
%======================================================================

\subsection{列并行数据流}

\begin{figure}[t]
\centering
\begin{tikzpicture}[font=\footnotesize, >=Stealth,
  mlp/.style={draw, fill=blue!10, minimum width=2.1cm, minimum height=7mm, align=center},
  ram/.style={draw, fill=yellow!15, minimum width=1.6cm, minimum height=7mm, align=center}]
\node[mlp, fill=gray!15] (f) at (0,0) {feeder MLP72};
\node[ram, left=6mm of f] (fb) {激活 BRAM\\（双缓冲）};
\node[mlp, above=6mm of f] (s1) {级 1 MLP72};
\node[mlp, above=6mm of s1] (s2) {级 2 MLP72};
\node[above=2mm of s2] (dots) {$\vdots$};
\node[mlp, above=2mm of dots] (sd) {级 $d$ MLP72};
\node[ram, right=6mm of s1] (w1) {列 1 权重\\BRAM72K};
\node[ram, right=6mm of s2] (w2) {列 2 权重\\BRAM72K};
\node[ram, right=6mm of sd] (wd) {列 $d$ 权重\\BRAM72K};
\draw[->] (fb) -- (f);
\draw[->, very thick, blue!60] (f) -- node[left, font=\scriptsize]{激活字} (s1);
\draw[->, very thick, blue!60] (s1) -- (s2);
\draw[->, very thick, blue!60] (s2) -- (dots);
\draw[->, very thick, blue!60] (dots) -- (sd);
\draw[->] (w1) -- (s1); \draw[->] (w2) -- (s2); \draw[->] (wd) -- (sd);
\node[right=18mm of w2, text width=3.6cm, align=left, font=\scriptsize] (note) {每级：同一激活字 $\times$ 本级一列权重，在本级累加器中累加本列列和；一次行传递产出 $d$ 个输出列。};
\end{tikzpicture}
\caption{列并行 MLP72 链的数据流：激活字沿硬级联向上广播，每一级持有一个输出列的权重并独立累加。}
\label{fig:chain}
\end{figure}

GEMM 链是阵列的核心计算单元（图~\ref{fig:chain}）。链由一个 feeder MLP72 与 $d$ 个乘法级 MLP72 组成（$d=16$ 或 32）。feeder 从自己的 BRAM 读出一个激活字，沿硬级联向上转发；每一级 MLP72 用自己 BRAM72K 中的一列输出权重与该字相乘，并在本级累加器中累加该列的部分和。一次 $W$ 个字的行传递产出 $d$ 个输出列；一个分块覆盖 $d\times\lfloor512/W\rfloor$ 列；行周期为
\begin{equation}
  T_{\text{row}} = \max\bigl(W+1,\; d,\; \sim\!14\bigr)\ \text{个阵列周期}.
  \label{eq:row}
\end{equation}

\textbf{为什么选择列并行。} MLP72 的级联只能把一个操作数字转发给同列的下一个 MLP72。若按常见的“部分和沿级联向上累加”组织，每一级只贡献点积的一段；按简报的分析，这种组织会浪费每级 16 个乘法器中的 15 个。列并行组织让同一个激活字在每一级与不同的列权重相乘，使每一级的 16 个（INT4 下 32 个）乘法器在一次行传递中全部忙碌。代价是每一级都需要自己的权重存储（BRAM72K）与结果引出通路，这是布线压力的主要来源（第~\ref{sec:impl} 节）。

\textbf{行周期的含义。} 由式~\eqref{eq:row}：当 $W$ 足够大（$W\ge d$ 且 $W\ge14$）时，行周期由 $W$ 决定，每个周期都有新字进入，乘法器接近满载；当 $W$ 较小时，行周期被链深 $d$ 或约 14 个周期的固定开销限制，乘法器出现空闲。另一方面，$W$ 越大，每级 BRAM 能容纳的列数 $\lfloor512/W\rfloor$ 越少，分块越多，权重装载次数越多。编译器的分块策略就是在这两端之间选取平衡点。

\subsection{分块与数据复用}

对一个 $M\times K\times N$ 的 GEMM，链上的执行可以理解为三重循环：外层遍历 $N$ 方向的分块（每块 $d\times\lfloor512/W\rfloor$ 列），中层遍历 $M$ 行，内层对每行做一次或多次行传递。在一个分块内，权重装入各级 BRAM 后被全部 $M$ 行复用，因此权重复用度等于 $M$；激活每行被 $d$ 个列复用，因此激活复用度等于 $d$。这一结构解释了两类负载的不同瓶颈：前缀 GEMM 的 $M\approx525$，权重装载可以被充分摊薄，瓶颈转为激活的反复重流（每个 $N$ 分块都要把全部 $M$ 行激活重读一遍）；动作专家 GEMM 的 $M=51$，每个分块的权重只被 51 行复用，装载成本难以摊薄。

为定量理解，本报告在两个假设下做一个示意推算 \tE：（a）一个激活字含 16 个 INT8 元素；（b）一个分块的权重沿链串行装入、每周期一个字，且不与计算重叠。此时一个分块的装载时间与计算时间之比约为 $d/M$：对 32 级链，前缀 GEMM 约 6\%，专家 GEMM 约 63\%；对 16 级链，二者约 3\% 与 31\%。这说明装入与计算的重叠对专家阶段至关重要，也解释了简报中“串行权重装载约占 15\%”\tP{} 的来源。编译器以双缓冲 feeder 和“权重装载提升到等待之前”缓解这一问题。

\subsection{feeder 与激活缓冲}

feeder 是链与外存之间的激活接口。它从本地 BRAM 读出激活字并送入级联，同时需要从外存预取下一行（或下一批行）的激活。简报所称的“双缓冲 feeder”，是指激活缓冲分为两半：一半供当前行传递使用，另一半同时从外存装入。双缓冲能否完全隐藏装入时间，取决于一行激活的装入时间是否短于一次行传递的时间；当 $W$ 较大时，行传递时间较长，装入容易被隐藏；当 $W$ 较小或外存访问受同步与竞争影响时，就会出现 feeder 等待，这对应时延构成中约 6\% 的“feeder 停顿”\tP。

\subsection{链深的选择}

链深 $d$ 决定了一次行传递产出的输出列数，也决定了行周期的下限。由式~\eqref{eq:row}，当 $W\ge d$ 时，链深不影响行周期，更深的链只是一次产出更多列、减少 $N$ 方向的分块数与激活重流次数；当 $W<d$ 时，行周期被 $d$ 限制，更深的链反而降低乘法器的有效利用率。对表~\ref{tab:gemm} 中的线性层，若一个激活字含 16 个 INT8 元素 \tE，则 $W$ 约为 $K/16$：视觉阶段约 72，前缀阶段约 128（down 投影分段后仍较大），动作专家阶段约 64 至 256。这些值都大于 32，因此对线性层而言 32 级链不会被行周期下限拖累。深度的限制主要来自另一端：更深的链需要更长的级联与更多的结果引出，布线与时序更难收敛。本设计同时提供 16 级与 32 级两种链，是在“每条链的产出”与“可布线的链数”之间的折中。

\subsection{累加宽度与溢出分析}

整数乘加的累加器必须保证在最长的归约维上不溢出。以最坏情况估计 \tE：
\begin{itemize}
  \item INT8$\times$INT8 线性层：单个乘积的绝对值不超过 $128\times128=16384$；最长的归约维为前缀 down 投影的 $K=16384$，按简报分为 4 段后每段 4096，每段累加和的绝对值不超过 $16384\times4096\approx6.7\times10^7$，远小于 int32 的 $2^{31}\approx2.1\times10^9$；即使不分段，$16384\times16384\approx2.7\times10^8$ 也在范围内。
  \item uint8$\times$int8 的 PV 运算：单个乘积绝对值不超过 $255\times128=32640$；归约维为键的个数（前缀 525，专家阶段 576），累加和绝对值不超过约 $1.9\times10^7$。
  \item INT4$\times$INT4：单个乘积绝对值不超过 64，累加更为宽裕。
\end{itemize}
因此，在当前数值方案下，int32 累加不存在溢出风险；K 分段的动机是片上缓冲容量与调度灵活性，而不是溢出。块浮点模式采用 fp24 累加，其误差来源是舍入而非溢出，已纳入位真模型。

\subsection{链的控制时序}

一条链的一次完整分块执行包括三个阶段：权重装载（把一个分块的列权重写入各级 BRAM72K）、行流过（逐行把激活字送入 feeder 并沿级联上传）、结果引出（把各级累加器的结果经 epilogue 写回外存）。三个阶段在不同分块之间可以流水：当前分块的行流过期间，可以装载下一个分块的权重（需要各级权重缓冲的双缓冲或分区）；当前分块的结果引出期间，可以开始下一个分块的行流过。简报列出的“双缓冲 feeder”与“权重装载提升到等待之前”属于这一类重叠。对 $M$ 较小的动作专家 GEMM，行流过阶段很短，权重装载的重叠是否充分直接决定链的有效利用率。

\subsection{K 分段与部分和合并}

$K$ 很大的 GEMM（典型如前缀的 down 投影，$K=16384$）需要沿 $K$ 方向分段。按简报，前缀 down 投影分为 4 段。每段在链上产出一组 int32 部分和，经反量化后以 fp32 形式合并。基线做法是由向量节点读回各段部分和并相加；K-split 反量化 epilogue 则让链自己从 GDDR6 读取前一段的部分和，在结果通路上完成反量化与合并。按简报，这一 epilogue 在 RTL 仿真中使对应向量阶段减少 81\%、整 chunk 减少 6.4\% \tMsim，比特流构建正在进行。

\subsection{结果通路 epilogue}

链的原始结果是每列的 int32 累加和。在基线设计中，这些 int32 结果写回 GDDR6，再由向量节点读回、反量化、施加激活函数，再写回。结果通路 epilogue 把这些步骤前移到链内：在结果离开节点之前完成反量化（行尺度 $\times$ 列尺度）与 GELU/GeGLU，并保证与向量节点上的同名运算逐位一致。其收益有两方面：省去一次 int32 结果的外存往返，以及省去一次独立的向量遍历。按简报，在同一比特流上对比，带 GeGLU/GELU epilogue 的 INT4 链使整 chunk 从 1.115~s 降到 1.047~s，即 $-6.1\%$ \tM，且结果与无 epilogue 版本逐位相同。

逐位一致是 epilogue 设计的硬约束：它要求 epilogue 中的查表、舍入与饱和行为与向量节点完全相同。本报告建议把两者的算子实现抽取为同一份参数化描述，由同一份位真软件模型验证，以降低维护两套实现的风险 \tR。

\subsection{PV 链}

注意力的“概率 $\times$ 值”运算的两个操作数都是运行时数据：概率矩阵由 softmax 产生并量化为 uint8，值矩阵为 int8。它与线性层 GEMM 的区别在于没有静态权重，且其 $K$ 维（键的个数，前缀为 525，专家阶段为 576）与 $N$ 维（头宽 256）的比例不同于线性层。设计上用 3 条独立的 16 级 PV 链（uint8 $\times$ int8）承担这部分运算，使注意力与线性层 GEMM 可以在不同节点上并行推进；编译器以注意力行块与层间重叠组织其调度。

\subsection{注意力数据流}

注意力是一个 chunk 中唯一两个操作数都为运行时数据的矩阵运算。其数据流可分为四步：$QK^\top$ 计算注意力分数；向量节点做 softmax 并把概率量化为 uint8；PV 链计算概率与值的乘积；结果经反量化进入输出投影。本设计的特点是把第三步交给专门的 PV 链：它的一个操作数（概率）是无符号的 uint8，与线性层的有符号 INT8 不同；它的归约维是键的个数（前缀 525、专家阶段 576），与头宽 256 的比例也不同于线性层。把 PV 运算放在独立的链上，一方面可以为 uint8$\times$int8 选择合适的乘法模式，另一方面可以让注意力与线性层在不同节点上并行，缩短关键路径。

“注意力行块”是编译器组织注意力的方式：把查询按若干行一组切分，每组的 softmax 与 PV 可以独立完成。这样，一组查询的 PV 可以与下一组查询的 softmax 重叠，也可以与下一层线性层的权重装载重叠。对动作专家而言，每步每层都要读取该层的前缀 KV；本报告建议评估把前缀 KV 常驻片上对 PV 链的收益（第~\ref{sec:risk} 节）\tR。

\subsection{PV 运算的形状分析}

PV 运算的形状与线性层不同：以前缀注意力为例，每个查询头的概率矩阵为 $525\times525$，值矩阵为 $525\times256$，即 $M=525$、$K=525$、$N=256$，共 8 个查询头共享同一个值矩阵（1 个 KV 头）；动作专家阶段为 $M=51$、$K=576$、$N=256$ \tE。与线性层相比，其 $N$ 较窄（256），而 8 个查询头共享同一值矩阵这一特点意味着：若把值矩阵作为“权重”装入链的各级，它可以被 8 个头 $\times$ $M$ 行复用，复用度相当高。这可能是 PV 链采用 16 级而非 32 级的一个原因：$N=256$ 在 16 级链上只需 16 个输出列分块，边角浪费更少。简报没有说明 PV 链的具体映射方式，以上分析仅供评审参考 \tE。

\subsection{链变体的取舍对比}

表~\ref{tab:variants} 从乘法吞吐、精度、布线与时钟四个维度对比各链变体，数据全部来自简报。

\begin{table}[h]
\centering
\caption{链变体对比（依据简报）。}
\label{tab:variants}
\small
\begin{tabularx}{\textwidth}{l Y Y Y}
\toprule
变体 & 改变了什么 & 已知收益 & 已知代价 \\
\midrule
INT8 基线 & — & 1.221~s，位真，精度可用 \tM & 布线限制阵列规模 \\
INT4（SIGNED 4$\times$4） & 每级乘法吞吐翻倍 & 1.007~s \tM & 精度不可用（$44.9^\circ$）\tM \\
块浮点 BFP8 & 共享指数 + fp24 累加 & 软件误差 0.70\% 对 0.97\% \tMsw & 仅 RTL 仿真 \\
GeGLU/GELU epilogue & 激活移入结果通路 & $-6.1\%$，逐位相同 \tM & 结果通路逻辑增加 \\
K-split epilogue & 部分和合并移入链 & 向量阶段 $-81\%$，整 chunk $-6.4\%$ \tMsim & 构建中 \\
fabric-light & 权重与结果改走硬级联 & 每级逻辑块 21.2 $\to$ 2.4；912 个乘法 MLP72 \tM & 时钟降至 310~MHz，1.588~s \tM \\
融合链对 & 两链共享激活行 & 读流量 $-11\%$ 至 $-23\%$ \tMsim & 时延收益约 0 \tP \\
链分裂 & 32 级在运行时拆为 2 $\times$ 16 & 比最佳固定形状好 0.02\% 至 0.12\% \tP & 需要第二个 feeder \\
\bottomrule
\end{tabularx}
\end{table}

从表中可以读出一个清晰的模式：改变乘法器本身（INT4）或阵列形状（融合、分裂）的变体，收益要么被精度否决，要么接近于零；减少链之外开销（epilogue、fabric-light 的布线节省）的变体，收益稳定但各自有限。这与校准模型“瓶颈在乘法器之外”的判断一致。

\subsection{链的变体}

\textbf{INT4 链。} 在 MLP72 的 SIGNED 4$\times$4 模式下，每级每周期 32 个有符号 INT4 乘加，乘法吞吐翻倍。在全部非注意力 GEMM 上使用 W4A4 时，整 chunk 为 1.007~s，但关节最大误差为 $44.9^\circ$ \tM，不可用。

\textbf{块浮点链。} 在 MLP72 的块浮点模式下使用 BFP8 操作数与 fp24 累加。软件模拟中其动作误差为 0.70\%，优于 INT8 的 0.97\% \tMsw；目前只有 RTL 仿真。

\textbf{fabric-light 链。} 权重经 BRAM72K 的写级联装入，结果经 MLP72 的输出级联分段引出，每级只需约 2.4 个逻辑块，而原方案约 21.2 个。该变体已在硅片上布下 912 个乘法 MLP72，整 chunk 为 1.588~s，结果位真 \tM；其时钟为 310 / 216.7 / 177.3~MHz，低于最佳 INT8 设计点。

\textbf{运行时融合链对。} 两条链经一条链路共享激活行，按 GEMM 在运行时选择作为一个两倍宽的阵列工作。RTL 仿真位真，读流量减少 11\% 至 23\% \tMsim，预计时延收益约为 0 \tP，比特流构建正在进行。

\textbf{链分裂（fission）。} 在级联中部静态配置第二个 feeder，使一条 32 级链可在运行时拆为两条 16 级链。该方案合法并已仿真；按校准模型，运行时切换只比最佳固定形状好 0.02\% 至 0.12\% \tP。

%======================================================================
\section{向量节点设计}\label{sec:vector}
%======================================================================

\subsection{功能}

8 个向量节点各有 2 条 lane，以 bf16 执行逐元素与归约运算。按简报，其职责包括：LayerNorm/RMSNorm、softmax、GeGLU、GELU、RoPE、残差相加、按 token 的量化（行尺度在运行时计算）、按“行尺度 $\times$ 列尺度”的反量化、融合的量化尾。超越函数（GELU、exp、rsqrt、sigmoid）基于查表。

\subsection{运算分类与数据通路}

从数据通路角度，向量运算可分为三类：
\begin{enumerate}
  \item \textbf{逐元素运算}：反量化、GELU/GeGLU、残差相加、RoPE。每个输出元素只依赖对应输入元素与少量参数，数据通路为“读—算—写”的流水线。
  \item \textbf{行归约运算}：RMSNorm/LayerNorm 的统计量、softmax 的最大值与指数和、按 token 量化的行尺度。需要先扫描整行得到统计量，再对整行做变换；数据通路需要行缓冲或两遍读取。
  \item \textbf{融合序列}：例如“归一化 $\to$ 量化”“反量化 $\to$ 激活 $\to$ 量化”。编译器把多个运算合并为一次遍历，减少外存往返；简报中的“融合量化尾”即属此类。
\end{enumerate}

\subsection{行归约运算的实现}

行归约运算（RMSNorm、LayerNorm、softmax、按 token 量化）需要先得到整行的统计量，再对整行做变换。在向量节点上有两种实现方式：一是\emph{两遍读取}，第一遍从外存读整行求统计量，第二遍再读整行做变换，外存流量加倍；二是\emph{行缓冲}，第一遍读入时把整行保存在片上缓冲中，第二遍从缓冲读取，外存流量不增加但需要片上存储。前缀阶段最宽的行（例如 GeGLU 之后的 16384 维中间激活需要按 token 量化）在 bf16 下约 32~KB/行 \tE，行缓冲是可行的。softmax 还可以采用在线算法，在一遍扫描中同时维护最大值与指数和，从而减少一次遍历。本报告建议在向量节点的设计评审中，对每一种行归约运算明确其实现方式及其外存流量，以便在校准模型中准确计入 \tR。

\subsection{向量运算的调度}

向量节点的工作在时间上分散于整个 chunk：每一层都有若干段向量运算，穿插在 GEMM 之间。其调度有两个特点。其一，向量运算的输入往往是多个 GEMM 分块的结果汇合，必须等全部分块完成才能开始，因此向量运算常常位于关键路径上；其二，向量节点只有 8 个，而同一时刻可能有多层、多种向量运算等待执行。编译器可以把同一层内彼此独立的向量运算（例如不同查询行块的 softmax）分配给不同的向量节点，并把向量运算与下一层 GEMM 的权重装载重叠。但当向量工作总量本身接近关键路径长度时，调度只能减少空闲，无法缩短必须完成的工作量；这时就需要结构性扩容或 epilogue 减负。

\subsection{查表超越函数}

超越函数以查表实现：输入的高位作为表地址，表项给出函数值或分段多项式系数。设计参数包括表的输入位宽、表项精度与插值方式。由于 epilogue 要求与向量节点逐位一致，这些参数一经确定就同时约束两处实现。本报告建议：以最终动作误差而非单个函数的误差为准选择表参数；把表内容作为编译产物的一部分纳入版本管理，并由位真模型生成 \tR。

\subsection{向量工作量的估算}

为理解向量通路为何成为瓶颈，本节按模型结构估算每个 chunk 的向量工作量 \tE。以“被向量节点读入或写出的元素个数”为粗略度量：
\begin{itemize}
  \item \textbf{前缀阶段}：每层至少包括两次 RMSNorm（输入宽 2048）、两次残差相加、RoPE（作用于 Q 与 K）、softmax（$8\times525\times525$ 个元素）、GeGLU（$2\times525\times16384$ 个输入元素），以及每个 GEMM 前后的量化与反量化。其中 GeGLU 与 MLP 相关的量化、反量化最为庞大，单层即有千万量级元素。
  \item \textbf{动作专家阶段}：每层的形状小一个数量级（51 个 token、宽 1024、MLP 4096），但要重复 180 次（18 层 $\times$ 10 步）；softmax 作用于 $8\times51\times576$ 个元素。
  \item \textbf{视觉阶段}：27 层，512 个 token，宽 1152，结构与前缀类似但规模较小。
\end{itemize}
以千万元素量级的单层工作量乘以数十至上百层，每个 chunk 的向量工作量在 $10^9$ 元素量级 \tE。16 条 lane 在约 290~MHz 下的理论吞吐约 $4.6\times10^9$ 元素/秒 \tE（假设每 lane 每周期处理一个元素）；若考虑行归约的两遍读取、外存往返与同步等待，实际有效吞吐会显著低于这一值。这一粗估与“向量工作约占 0.61~s”\tD{} 的量级相容，说明向量通路的瓶颈是结构性的，而不是某个算子实现不佳。上述估算依赖本报告的假设（每 lane 每周期一个元素），真实的 lane 宽度与每周期吞吐以设计文件为准。

\subsection{epilogue 与向量节点的分工原则}

哪些运算应当留在向量节点、哪些应当并入链的结果通路，可以按三条原则判断：
\begin{enumerate}
  \item \textbf{局部性}：只依赖本元素与少量常数的运算（反量化、激活函数、残差相加）最适合并入结果通路；需要整行统计量的运算（归一化、softmax、按 token 量化）并入结果通路需要跨列归约，代价更高。
  \item \textbf{流量}：若某运算的输入是链的 int32 结果，把它并入结果通路可以同时省去一次 int32 写回与一次读回，收益最大。
  \item \textbf{一致性}：并入结果通路的运算必须与向量节点逐位一致；若某运算的查表或舍入难以在结果通路上复制，就应留在向量节点。
\end{enumerate}
已上板的 GeGLU/GELU epilogue 与已仿真的 K-split 反量化 epilogue 都符合这三条原则。按这一原则，下一批候选是反量化 + 残差相加，以及按 token 量化的行尺度计算（需要结果通路具备行归约能力）。

\subsection{向量通路是当前瓶颈}

按简报的校准模型，最佳 INT8 设计点的 1.221~s 中向量工作约占 50\% \tP；INT4 之后，约 82\% 的关键路径是向量工作 \tP。这两个占比口径不同（前者为时间构成，后者为关键路径占比），时钟也不同，本报告不把它们相乘换算为绝对时间。但结论是一致的：向量通路的吞吐没有随矩阵阵列同比例扩展。16 条 lane（8 节点 $\times$ 2 lane）对应的是一个在每两个 GEMM 之间都要执行数次完整遍历的负载。第~\ref{sec:risk} 节把向量通路的结构性扩容列为最高优先级的后续工作。

\subsection{向量通路扩容方案比较}

扩大向量通路吞吐有四种途径，各有不同的代价 \tR：
\begin{itemize}
  \item \textbf{增加向量节点数}：最直接，但每个节点都带有程序获取、外存访问与同步的固定开销，并占用一个 NAP；在布线已经受限的情况下，新增节点可能迫使减少链节点。
  \item \textbf{加宽每个节点的 lane}：分摊了节点的固定开销，但更宽的 lane 需要更宽的外存访问与行缓冲，对单个节点的布线局部性要求更高。
  \item \textbf{提高向量域时钟}：当前向量域为 248.4 至 288.9~MHz，低于阵列域；若能通过流水化与重定时提高频率，收益与资源无关。但浮点与查表逻辑的时序本身较紧，提频空间需要实际评估。
  \item \textbf{降低向量运算的位宽}：在精度允许的前提下，把部分 bf16 运算改为更窄的定点或块浮点，可以在同样的资源下提高吞吐；但需要以动作误差为准重新验证，且会改变 epilogue 与向量节点逐位一致的约束。
\end{itemize}
此外，把更多运算移入链的结果通路（epilogue 扩展）是一种“不扩容而减负”的途径。上述途径可以组合；本报告建议先用校准模型量化每种途径对总时延的边际收益，再结合资源预算选择组合。

%======================================================================
\section{数值方案}\label{sec:numerics}
%======================================================================

\subsection{精确路径（W8A8）}

按简报，精确路径的数值方案为：
\begin{itemize}
  \item \textbf{权重}：每个 GEMM 的权重按输出通道量化为 INT8，截断阈值按均方误差（MSE）选取；
  \item \textbf{激活}：按 token 动态量化为 INT8，行尺度在运行时计算；
  \item \textbf{语言模型}：采用 SmoothQuant~\cite{smoothquant}，把激活中的离群值通过等价变换迁移到权重；
  \item \textbf{注意力概率}：uint8；
  \item \textbf{向量接口}：所有向量运算的输入输出为 bf16。
\end{itemize}
在这一方案下，46 帧硅片实测的关节最大误差为 $0.665^\circ$、平均误差为 $0.072^\circ$，相对 fp32 参考 \tM。

\subsection{数据流中的数值转换}

一个线性层在精确路径上的数值转换顺序如下：上游向量运算输出 bf16 激活 $\to$ 按 token 求行尺度并量化为 INT8 $\to$ 链上做 INT8$\times$INT8 乘加，累加为 int32 $\to$ 以“行尺度 $\times$ 列尺度”反量化 $\to$ 以 bf16 进入下游向量运算。K 分段的 GEMM 在反量化后以 fp32 合并部分和。这一顺序的每一步都在位真软件模型中精确建模，舍入模式与饱和行为与硬件一致，这是逐位核对的前提。

\subsection{SmoothQuant 的实现位置}

SmoothQuant~\cite{smoothquant} 的核心是一个数学等价的逐通道缩放：把激活的第 $j$ 个通道除以 $s_j$，同时把权重的第 $j$ 行乘以 $s_j$，使激活中的离群通道被“压平”，量化难度转移到更易量化的权重上。在本设计中，这一缩放可以完全离线完成：权重侧的乘法在量化前并入权重，激活侧的除法并入上游的归一化参数（例如 RMSNorm 的逐通道增益），因而在运行时不增加任何额外运算。这是 SmoothQuant 适合硬件实现的重要原因；需要确认的只是位真模型与编译器对缩放后的参数使用一致。

\subsection{注意力概率的 uint8 表示}

softmax 的输出在 $[0,1]$ 之间，以 uint8 表示时量化步长为 $1/255$。这一选择的好处是 PV 运算可以在整数乘法器上完成，且无符号表示比有符号表示多一位有效精度；代价是小于约 $1/510$ 的概率被量化为 0。对于前缀 525 个键的注意力，大多数键的概率可能很小，被量化为 0 的比例可能不低；这对精度的影响已经包含在 46 帧的实测误差中。若后续的敏感度分析显示注意力是主要误差来源，可以考虑在概率量化中引入逐行尺度 \tR。

\subsection{bf16 接口的取舍}

所有向量运算以 bf16 为接口。bf16 的动态范围与 fp32 相同，适合承载归一化前的激活与残差流，避免溢出；其 8 位尾数带来约 0.4\% 的相对舍入误差。相比 fp32 接口，bf16 使向量节点与外存之间的流量减半；相比 INT8 接口，它避免了在每个向量运算前后都做量化。在向量通路成为瓶颈的情况下，接口位宽直接影响向量工作的外存流量；是否在部分接口上采用更窄的格式，应作为敏感度分析的一项内容 \tR。

\subsection{低比特与块浮点}

\textbf{INT4。} W4A4 在全部非注意力 GEMM 上的结果不可用（关节最大误差 $44.9^\circ$ \tM）。简报的判断是“4 位权重对该模型是破坏性的”。这一结果只说明均匀 W4A4 不可用，并不排除逐层混合精度；任何混合方案都必须以动作误差为准重新验证。

\textbf{BFP8。} 块浮点让一组元素共享指数，元素本身用窄位宽尾数~\cite{msfp,mx}。MLP72 支持块浮点模式；以 BFP8 操作数与 fp24 累加时，软件模拟的动作误差 0.70\% 优于 INT8 的 0.97\% \tMsw。块浮点还可能简化尺度处理，因为块内共享指数使部分反量化工作可以在链内完成；这一点需要实际设计确认。

\subsection{误差来源分析}

精确路径相对 fp32 参考的误差来自以下几个环节：
\begin{enumerate}
  \item \textbf{权重量化}：按输出通道的 INT8 量化与 MSE 截断。截断阈值越小，量化步长越细，但被截断的离群权重误差越大。
  \item \textbf{激活量化}：按 token 的动态 INT8 量化。每行一个尺度，行内离群值会压缩其他元素的有效位数；SmoothQuant 通过把离群值迁移到权重来缓解这一问题。
  \item \textbf{注意力概率量化}：softmax 输出以 uint8 表示，接近 0 的小概率被量化为 0，接近 1 的大概率的分辨率有限。
  \item \textbf{bf16 接口}：所有向量运算的输入输出为 bf16，其 8 位尾数带来约 $2^{-8}$ 的相对误差。
  \item \textbf{查表近似}：超越函数的查表与插值误差。
  \item \textbf{迭代累积}：动作专家的 10 个积分步中，每一步的误差都会传入下一步。
\end{enumerate}
这些误差在最终动作上的合成效应，就是 46 帧上的关节最大误差 $0.665^\circ$ 与平均误差 $0.072^\circ$ \tM。INT4 的失败（$44.9^\circ$ \tM）说明，把第 1、2 项的位宽再减半会使误差越过可用边界。

\subsection{敏感度分析方法}

为给后续的混合精度或块浮点方案提供依据，本报告建议以下敏感度分析流程 \tR：
\begin{enumerate}
  \item \textbf{逐层替换}：在位真模型中，每次只把一层（或一类 GEMM）改为候选格式（INT4、BFP8 等），其余保持 INT8，测量 46 帧上的动作误差；
  \item \textbf{按阶段替换}：分别只替换视觉、前缀或动作专家阶段，判断误差主要来自哪个阶段；
  \item \textbf{按环节替换}：分别只改变权重、激活或注意力概率的格式，判断误差主要来自哪个环节；
  \item \textbf{组合验证}：根据前三步的结果选出低敏感度的层组合，再整体验证；
  \item \textbf{闭环确认}：对候选方案做闭环评测，确认开环误差的改善转化为任务表现。
\end{enumerate}
由于位真模型与硬件逐位一致，上述分析得到的精度结论可以直接迁移到硅片上，不需要每个候选都上板。

\subsection{块浮点的实现要点}

在 MLP72 的块浮点模式下实现 BFP8 链，需要处理以下问题 \tR：
\begin{itemize}
  \item \textbf{块的划分}：共享指数的块沿 $K$ 方向划分，使同一块内的元素在同一个点积中相乘累加；块大小需要与激活字的宽度和链的行传递对齐。
  \item \textbf{指数对齐}：块内元素以共享指数对齐尾数，乘加在尾数域完成，块与块之间的累加需要处理不同的共享指数，这由 fp24 累加器承担。
  \item \textbf{与向量节点的接口}：若激活以 BFP8 格式进入链，则上游向量节点需要输出 BFP8（计算每块的共享指数），这替代了当前按 token 求行尺度的 INT8 量化；反量化也相应变为以块指数恢复数值。
  \item \textbf{位真建模}：块浮点的舍入与对齐规则必须在位真模型中精确复现；软件模拟中的 0.70\% 误差 \tMsw{} 是否与硬件逐位一致，需要在 RTL 仿真与上板时确认。
\end{itemize}
块浮点可能带来的一个附带收益，是把“按 token 量化”的行归约替换为“按块求指数”的局部运算，减少向量节点的行归约工作量；这一点需要在具体设计中量化。

\subsection{位真软件模型}

位真软件模型是数值方案与验证体系之间的桥梁。它以与硬件相同的整数宽度、舍入与饱和规则，实现所有 GEMM、epilogue 与向量运算；其输出既用于评估精度（与 fp32 参考比较），也被编译器用于生成期望内存（第~\ref{sec:compiler} 节）。因此，一次“精度评估”与一次“硬件正确性验证”使用的是同一份数值定义：若硅片结果与期望内存逐位一致，则硅片的精度就等于位真模型的精度。这一性质使精度问题与正确性问题可以被干净地分离：精度不达标是数值方案的问题，结果不一致是硬件或编译器的问题。

%======================================================================
\section{编译器：chunk 生成器}\label{sec:compiler}
%======================================================================

\subsection{职责}

chunk 生成器把整个 chunk 降低为 35 个节点各自的程序，并同时生成 GDDR6 的初始内容（程序、权重、常数表）与每一步的期望内存。其输入是模型结构与量化后的参数、目标设计点的硬件描述（链的数量与深度、向量节点数、时钟、可用的 epilogue 与运行时模式）；输出是可直接装入 GDDR6 的镜像与用于核对的期望内存。

\subsection{编译流程}

\begin{figure}[t]
\centering
\begin{tikzpicture}[font=\footnotesize, >=Stealth, node distance=4mm,
  st/.style={draw, rounded corners=2pt, fill=#1, text width=2.5cm, align=center, minimum height=9mm},
  st/.default=blue!8]
\node[st=gray!10] (in) {模型结构 + 量化参数 + 硬件描述};
\node[st, right=of in] (lower) {算子图降低\\GEMM / 注意力 / 向量};
\node[st, right=of lower] (tile) {按深度类分块\\分配节点};
\node[st, right=of tile] (mem) {GDDR6 地址映射\\（条带化）};
\node[st, below=8mm of mem] (sched) {调度：重叠、提升、融合、剪除};
\node[st, left=of sched] (sync) {插入 WAIT/POST\\负载均衡};
\node[st, left=of sync] (emit) {生成节点程序\\与 GDDR6 镜像};
\node[st=yellow!15, left=of emit] (exp) {位真模型\\生成期望内存};
\draw[->] (in) -- (lower); \draw[->] (lower) -- (tile); \draw[->] (tile) -- (mem);
\draw[->] (mem) -- (sched); \draw[->] (sched) -- (sync); \draw[->] (sync) -- (emit); \draw[->] (emit) -- (exp);
\end{tikzpicture}
\caption{chunk 生成器的编译流程。}
\label{fig:compiler}
\end{figure}

编译流程如图~\ref{fig:compiler} 所示，按简报可归纳为以下步骤。

\textbf{（1）算子图降低。} 把模型展开为一个 chunk 内的完整算子序列：27 层 SigLIP、18 层 Gemma-2B 前缀、18 层 $\times$ 10 步的动作专家，以及其间的全部注意力与向量运算。

\textbf{（2）按深度类分块。} 对每个 GEMM，根据其 $K$、$N$ 与 $M$ 选择链的深度类（16 级或 32 级）、每行传递的字数 $W$ 与分块大小，使式~\eqref{eq:row} 的行周期接近 $W+1$，同时控制分块数量。

\textbf{（3）条带化地址映射。} 把权重、激活与 KV 按块分散到多个 GDDR6 通道上，使多个节点的并发访问均匀分布，避免通道热点。

\textbf{（4）调度优化。} 包括：双缓冲 feeder，使下一行或下一分块的激活装入与当前计算重叠；把权重装载提升到 \texttt{WAIT} 之前，使等待上游结果期间可以预取权重；融合量化，把量化并入上游向量运算的尾部；注意力行块与层间重叠；剪除无用工作，例如前缀最后一层只计算其 K/V。

\textbf{（5）同步插入与负载均衡。} 为每一对生产者—消费者插入 \texttt{POST}/\texttt{WAIT}；把 GEMM 的分块在同类节点之间分配，使各节点的完成时间接近。

\textbf{（6）程序与镜像生成。} 生成每个节点的程序与 GDDR6 初始镜像。

\textbf{（7）期望内存生成。} 以位真软件模型执行同一调度，生成每一个写入 GDDR6 的数据在完成时刻的期望值。简报称这一“期望内存模型”使每一次硅片运行都可以逐拍核对。

\subsection{调度问题的形式}

从编译器的角度，chunk 调度可以表述为一个带资源约束的任务调度问题：任务是 GEMM 分块、注意力行块与向量运算序列；资源是 24 个 GEMM 链节点（分两个深度类）、3 个 PV 链节点与 8 个向量节点；约束包括数据依赖、每类任务只能分配给相应类型的节点，以及外存带宽。目标是最小化整个 chunk 的完成时间。由于调度完全静态，编译器可以在编译期对整个 chunk 做全局优化，而不需要运行时决策。

简报列出的优化手段可以放在这个框架中理解：
\begin{itemize}
  \item \textbf{负载均衡}：把一个 GEMM 的分块在同类节点之间分配，使它们的完成时间接近，减少关键路径上的空闲；
  \item \textbf{重叠}：注意力行块与层间重叠，使 PV 链、GEMM 链与向量节点在不同层的工作上并行；
  \item \textbf{提升}：把不依赖上游结果的权重装载移到 \texttt{WAIT} 之前，缩短等待期间的空闲；
  \item \textbf{融合与剪除}：减少任务数量与外存往返。
\end{itemize}

\subsection{分块参数的选择}

对每个 GEMM，编译器需要确定三个参数：深度类（16 级或 32 级链）、每次行传递的字数 $W$（对应 $K$ 方向的切分），以及 $N$ 方向每个分块覆盖的列数 $d\times\lfloor512/W\rfloor$。选择的依据是一个简单的代价模型：
\begin{itemize}
  \item \textbf{计算代价}：分块数 $\times$ $M$ $\times$ 每行传递次数 $\times$ $T_{\text{row}}$；
  \item \textbf{装载代价}：分块数 $\times$ 每分块的权重字数（若不能与计算重叠）；
  \item \textbf{激活重流代价}：$N$ 方向分块数 $\times$ $M$ 行激活的外存读取量；
  \item \textbf{结果代价}：输出元素数 $\times$ 每元素的结果引出与写回代价（若有 epilogue，则部分代价转移到结果通路）。
\end{itemize}
对前缀的宽 GEMM，激活重流代价占主导，编译器倾向于选择更深的链以减少 $N$ 方向的分块数；对动作专家的 GEMM，装载代价占主导，编译器倾向于让装载与计算充分重叠，并在多个节点之间分摊一个 GEMM 的分块，使每个节点装载的权重更少。这一代价模型与第~\ref{sec:impl} 节的资源预算、第~\ref{sec:verif} 节的校准模型共用同一组硬件参数。

\subsection{负载均衡}

同一个 GEMM 的分块可以分配给同类的多个节点。负载均衡的目标是使这些节点的完成时间尽可能接近，因为下游运算要等待全部分块完成。简单的均分在两种情况下会失衡：一是分块数不能被节点数整除，最后一轮只有部分节点工作；二是边角分块（$N$ 方向不满的最后一块）比其他分块小。编译器可以把不同 GEMM 的分块交错分配，使一个 GEMM 的尾部与下一个 GEMM 的头部重叠，从而填平尾部空闲；这要求两个 GEMM 之间没有数据依赖，或者依赖已被满足。对 16 级与 32 级两类链，编译器还可以把同一个 GEMM 分配给两类链的混合，以更细的粒度平衡完成时间。

\subsection{期望内存的生成}

期望内存的生成方式是：编译器按其生成的程序顺序，以位真模型逐条“执行”每个节点的命令，记录每一次外存写入的地址与数据。由于命令流不含数据依赖的分支，这一执行与硬件执行在每一个外存写入上都应一致；节点之间的相对时序不影响结果，因为所有跨节点的数据依赖都由 \texttt{WAIT}/\texttt{POST} 显式保护。这也提出了一个必须满足的正确性条件：\emph{编译器插入的同步必须覆盖所有真实的数据依赖}。若遗漏一个依赖，硬件可能在数据就绪前读取，得到与期望内存不同的结果；逐拍核对正好可以暴露这类错误。

\subsection{编译产物与接口}

编译器为每个设计点、每个模型检查点输出一组产物：各节点的命令流、GDDR6 的初始镜像（程序、权重、常数与查表内容）、输入与输出区域的地址约定、期望内存，以及供性能模型使用的静态统计（每个节点的命令数、分块数与预计周期）。主机程序只依赖输入与输出区域的地址约定和启动/完成标志的位置，而这些约定在不同设计点之间保持不变，从而满足需求 N4。期望内存体积很大，验证时通常只比对输出区域与抽样的中间区域；在调试时再比对全部区域。

\subsection{编译器自身的验证}

编译器是验证链中的关键环节，它的错误会同时影响程序与期望内存，因而可能出现“两边一起错、核对却通过”的情况。为避免这一风险，本报告建议 \tR：
\begin{itemize}
  \item 以独立于编译器的 L1 位真模型（直接按模型结构计算，而不按调度执行）作为期望内存的对照，确认调度没有改变数值结果；
  \item 对同步插入做静态检查：遍历所有跨节点的读写对，确认每一对之间都有 \texttt{POST}/\texttt{WAIT} 保护；
  \item 对地址映射做静态检查：确认不同张量的地址区间在其生命周期内不重叠；
  \item 以随机调度扰动（在合法范围内改变节点分配与顺序）运行 L3 仿真，确认结果不变。
\end{itemize}

\subsection{调度优化的效果}

在同一套 INT8 硬件上，调度优化把整 chunk 时延从 1.901~s 降到 1.221~s \tM，约 1.56 倍 \tD（$1.901/1.221$）。简报列出的手段依次为：内存布局、排序与重叠、融合、死代码剪除，最后是时钟。第一次整阵列运行时延为 6.48~s \tM，到最佳 INT8 点累计改善约 5.3 倍 \tD（$6.48/1.221$），其中也包含时钟与硬件迭代的贡献。这组数据说明，在本架构中编译器与调度是与硬件同等重要的性能杠杆。

\subsection{调度的局限与改进空间}

当前的静态调度已经带来约 1.56 倍的改善 \tD，但仍有以下局限 \tR：
\begin{itemize}
  \item \textbf{步间依赖不可打破}：动作专家的 10 个积分步严格串行，调度只能在一步之内做重叠；
  \item \textbf{向量工作的总量不变}：调度只能消除空闲，无法减少向量运算本身；
  \item \textbf{同步往返的固定开销}：每次 \texttt{WAIT}/\texttt{POST} 都要经过外存，调度只能通过减少同步次数（合并依赖）来降低其影响；
  \item \textbf{模型精度的局限}：调度依赖代价模型选择分块与分配，代价模型与硬件的偏差会导致次优选择。
\end{itemize}
改进空间主要在两方面：以校准后的模型对更大的调度空间做搜索（例如分块参数的联合优化）；以及在硬件支持片上驻留后，把前缀 KV 与部分专家权重的放置纳入调度决策。

\subsection{实验管理}

由于同一编译器服务多个设计点与多种调度选项，实验管理本身成为一项工程工作 \tR。建议编译器的每一次运行都记录其输入配置（设计点、调度选项、模型检查点、量化参数）与输出摘要（预计周期、分块统计、期望内存的校验和），并与硅片或仿真的结果关联。这样，任何一个时延或精度数字都可以追溯到产生它的完整配置；设计空间探索中“哪一项改动带来了哪一部分收益”也可以被清楚地回答。简报中“调度带来 1.901~s 到 1.221~s”的逐项归因（内存布局、排序与重叠、融合、死代码剪除、时钟）正是这类实验管理的产物。

\subsection{设计点无关性}

同一编译器为多个设计点生成程序：INT8 链、INT4 链、带 epilogue 的链、fabric-light 链、融合链对与链分裂。设计点之间的差别体现为硬件描述中的参数，编译器据此选择可用的指令与调度策略。这使设计空间探索的迭代成本主要落在综合与布线上，而不是软件改写上。

%======================================================================
\section{FPGA 物理实现}\label{sec:impl}
%======================================================================

\subsection{实现流程}

设计以 Synplify 综合，以 Achronix ACE 10.5.2 完成布局、布线与时序分析。MLP72 的级联连接与乘法模式、BRAM72K 与 LRAM2K 的配对关系在综合前由参数确定；链的长度与位置需要与硬核的列结构对齐，因为硬级联只在同一列相邻的 MLP72 之间存在。

\subsection{资源利用与布线约束}

按简报，布线资源而非乘法器是阵列规模的约束：一个完整阵列使用 2,560 个 MLP72 中的约 760 至 860 个，逻辑块占用率为 63\% 至 75\%，更大的阵列无法完成布线。这一现象可以从列并行链的结构理解：每一级 MLP72 都需要独立的权重写入路径、控制信号与结果引出通路，这些都要经过通用 fabric 布线；链越多，这些分布式连线越多，最终先于乘法器耗尽。

fabric-light 链正是针对这一约束的结构性改进：把权重写入改由 BRAM72K 的写级联完成，把结果引出改由 MLP72 的输出级联分段完成，使每级所需的逻辑块从约 21.2 个降到约 2.4 个，并在硅片上布下 912 个乘法 MLP72 \tM。但该设计点的阵列时钟为 310~MHz，低于最佳 INT8 点的 500~MHz，整 chunk 为 1.588~s \tM。因此，“能布下更多乘法器”尚未转化为“更快”；原因在于时钟下降，以及向量通路没有同步扩展。

\subsection{资源利用现状的解读}

“约 760 至 860 个 MLP72、63\% 至 75\% 的逻辑块”\tM{} 这一资源状态值得从两个角度解读。从乘法器角度看，约三分之二的 MLP72 没有被使用，这是一种明显的浪费，说明器件的硬核资源远未被充分利用。从逻辑块角度看，63\% 至 75\% 的占用率在 FPGA 设计中已经偏高，继续增加逻辑往往会导致布线拥塞与时序恶化。两者合起来说明：限制阵列规模的不是“逻辑块不够”，而是“布线不够”；而布线消耗的主要来源是每级 MLP72 周围的分布式连线。fabric-light 链把这些连线大部分改由硬级联承担，是解决这一问题的正确方向。

\subsection{布局规划}

MLP72 与 BRAM72K 在器件上按列排布，硬级联只在同一列相邻的 MLP72 之间存在。因此，一条链在物理上必须占据同一列中连续的 MLP72，其 feeder 位于链底。布局规划的主要问题是：把 24 条 GEMM 链与 3 条 PV 链放到哪些列、哪些行；把 8 个向量节点放在哪里；以及每个节点的 NAP 位于何处。本报告建议的布局原则如下 \tR：
\begin{itemize}
  \item \textbf{链靠近其 NAP}：每条链的权重写入与结果引出都经 NAP 与外存交换数据，链与 NAP 之间的距离直接影响布线长度与时序；
  \item \textbf{向量节点分散放置}：向量节点与多条链交换数据（经外存），分散放置可以避免局部布线拥塞；
  \item \textbf{预留布线通道}：在链之间预留布线通道，给权重写入与结果引出的分布式连线留出空间；
  \item \textbf{按时钟域分区}：阵列域、向量域与 fabric 域的逻辑尽量分区放置，减少跨域连线的长度。
\end{itemize}
fabric-light 链把权重写入与结果引出改由硬级联完成，本质上就是减少了第一、三条所针对的分布式连线，因此能在同样的布线资源下放下更多链。

\subsection{资源预算方法}

每个设计点的资源消耗可以按节点类型分解：GEMM 链节点的 MLP72 数等于链深加 feeder，并各自占用配对的 LRAM2K 与若干 BRAM72K；逻辑块主要消耗在每级的权重写入、结果引出与控制上（原方案每级约 21.2 个，fabric-light 约 2.4 个，据简报）；向量节点的逻辑块消耗在浮点单元、查表与行缓冲上；每个节点还有一份程序获取、外存访问与同步逻辑。本报告建议在设计空间探索之前，先建立一个按节点类型的资源预算表，并用已实现的设计点校准，使“增加一条链”或“增加一个向量节点”的资源代价可以在综合前估算 \tR。这对第~\ref{sec:risk} 节的“矩阵/向量算力再配比”尤为重要。

\subsection{构建流程与版本管理}

一个设计点的完整构建包括：RTL 参数化（链数、链深、模式、epilogue 类型、时钟）、综合、布局布线、时序签核与比特流生成；同时，编译器需要该设计点的硬件描述以生成程序。本报告建议把“比特流版本、硬件描述、编译器版本、程序镜像、期望内存”作为一个整体进行版本管理 \tR：任何一次硅片运行都应能追溯到这五者的具体版本，以保证结果可复现。简报中的各设计点（最佳 INT8、INT4、INT4 + epilogue、fabric-light）应分别对应这样一组可追溯的构建产物。

\subsection{时钟与时序收敛}

表~\ref{tab:clocks} 汇总各设计点的时钟。

\begin{table}[h]
\centering
\caption{各设计点的时钟（阵列 / 向量 / fabric，MHz），依据简报。}
\label{tab:clocks}
\small
\begin{tabular}{l c}
\toprule
设计点 & 时钟 \\
\midrule
INT8，最佳调度 & 500 / 288.9 / 216.7 \\
INT4 链 & 437.5 / 279.6 / 244.6 \\
INT4 + GeGLU/GELU epilogue（对比用比特流） & 481.25 / 248.4 / 180.7 \\
fabric-light 链 & 310 / 216.7 / 177.3 \\
\bottomrule
\end{tabular}
\end{table}

从表中可以看出，时钟随设计点变化很大，且三个域的变化方向并不一致。这对设计评估有两个含义：其一，跨设计点比较时延时，时钟与结构的贡献混在一起，只有在同一比特流上的开关对比（如 epilogue 的 $-6.1\%$）才能给出干净的结构收益；其二，向量域时钟（248.4 至 288.9~MHz）明显低于阵列域，而向量工作又是第一瓶颈，因此提高向量域时钟本身就是一个直接的优化方向。

本报告建议的时序收敛措施包括 \tR：对链与其所在硬核列做区域约束，使每条链的权重写入与结果引出通路局部化；在跨长距离的控制与数据通路上增加流水级；对向量节点的查表与浮点单元做寄存器重定时；把 NAP 与节点的相对位置纳入布局规划，缩短节点到 NoC 的连线。

\subsection{功耗与热设计考虑}

简报没有给出功耗数据，本节只说明功耗评估需要考虑的因素 \tR。FPGA 板卡的功耗由静态功耗与动态功耗组成：动态功耗与各时钟域的频率、翻转率和使用的资源数量成正比。在本设计中，阵列域运行在最高频率，MLP72 与 BRAM72K 的翻转是主要的动态功耗来源；GDDR6 接口与 NoC 在持续的大流量访问下同样消耗可观的功率。由于调度静态，每个 chunk 的活动模式是确定的，因此每 chunk 能耗应当具有很小的波动，这使能耗测量本身可以作为确定性的旁证。评估机器人部署时，除了平均功耗，还应测量峰值功耗与长时间运行下的温度，以确认板卡在机器人整机的散热条件下能维持设计时钟。

\subsection{fabric-light 时钟下降的分析}

fabric-light 链在布线上的改进是显著的（每级逻辑块从约 21.2 降到约 2.4），但其阵列时钟从最佳 INT8 点的 500~MHz 降到 310~MHz，向量域与 fabric 域也分别降到 216.7 与 177.3~MHz \tM。可能的原因包括：更长的硬级联链路（权重写级联与结果输出级联）引入了更长的组合路径；更多的链带来更高的整体资源占用，使剩余逻辑的布线更加拥挤；以及结果分段引出所需的控制逻辑。简报没有给出时序报告，本报告无法确定具体原因，建议在该设计点的时序报告中重点检查：级联相关路径是否成为关键路径；向量域与 fabric 域的时钟下降是否源于整体拥塞；以及在级联的分段处增加寄存器是否可以恢复频率 \tR。由于时延约与时钟成反比，若能在保持 912 个乘法 MLP72 的同时恢复阵列与向量时钟，并同步扩容向量通路，fabric-light 有望成为后续的基线结构。

\subsection{时序签核与多频点}

由于同一套 RTL 在不同设计点上以不同的时钟组合运行（表~\ref{tab:clocks}），时序签核需要针对每一个发布的比特流单独完成，并记录三个时钟域各自的最差裕量与关键路径类别。建议在签核报告中把关键路径按类别统计（级联相关、NoC 接口、向量浮点、查表、跨域同步），以便在不同设计点之间比较，判断哪类路径最限制频率 \tR。这一统计也能为第~\ref{sec:impl} 节 fabric-light 时钟下降的分析提供直接证据。

\subsection{跨时钟域}

三个时钟域之间的数据经异步 FIFO 传递，控制信号经同步器传递。由于简报中的设计点时钟各不相同，且同一 RTL 需要在多个频点组合下工作，跨域逻辑必须对频率比不作任何假设。验证要求见第~\ref{sec:verif} 节。

%======================================================================
\section{验证方案与验证结果}\label{sec:verif}
%======================================================================

\subsection{验证目标}

验证体系要回答三个问题：
\begin{enumerate}
  \item 硬件与编译器是否正确实现了数值方案？（功能正确性）
  \item 数值方案是否满足精度需求？（精度）
  \item 性能模型是否准确预测了硬件行为？（性能可预测性）
\end{enumerate}
第一个问题的判据是逐位一致；第二个问题的判据是相对 fp32 参考的动作误差；第三个问题的判据是校准模型预测与实测之间的偏差。

\subsection{五级验证链}

\begin{figure}[t]
\centering
\begin{tikzpicture}[font=\footnotesize, >=Stealth, node distance=5mm,
  lv/.style={draw, rounded corners=2pt, fill=#1, text width=2.3cm, align=center, minimum height=11mm},
  lv/.default=blue!8]
\node[lv=gray!10] (l0) {L0 fp32 参考\\PyTorch};
\node[lv, right=of l0] (l1) {L1 量化位真\\软件模型};
\node[lv, right=of l1] (l2) {L2 编译器\\期望内存};
\node[lv, right=of l2] (l3) {L3 周期精确\\RTL 仿真（Verilator）};
\node[lv=green!10, right=of l3] (l4) {L4 硅片\\逐拍核对};
\draw[->] (l0) -- node[below, font=\scriptsize, text width=1cm, align=center, yshift=-5mm]{精度} (l1);
\draw[->] (l1) -- node[below, font=\scriptsize, yshift=-5mm]{调度} (l2);
\draw[->] (l2) -- node[below, font=\scriptsize, yshift=-5mm]{逐位} (l3);
\draw[->] (l3) -- node[below, font=\scriptsize, yshift=-5mm]{逐位} (l4);
\end{tikzpicture}
\caption{五级验证链。L0 与 L1 之间比较精度，L1 至 L4 之间要求逐位一致。}
\label{fig:vchain}
\end{figure}

验证体系由五级组成（图~\ref{fig:vchain}）：

\textbf{L0　fp32 参考。} 以 fp32 PyTorch 实现的 \pizero{} 作为精度基准。

\textbf{L1　量化位真软件模型。} 以与硬件相同的整数宽度、舍入与饱和规则实现全部运算。L0 与 L1 之间比较精度：46 帧录制数据上的关节角最大与平均误差。

\textbf{L2　编译器期望内存。} 编译器以 L1 的数值定义执行其生成的调度，得到每个外存写入的期望值。L1 与 L2 之间要求逐位一致，用于验证调度（分块、融合、剪除、K 分段合并顺序）没有改变数值结果。需要注意，K 分段的合并顺序会影响浮点累加的结果，因此 L1 必须采用与调度相同的合并顺序，或者把合并顺序作为 L1 的参数。

\textbf{L3　周期精确 RTL 仿真。} 以 Verilator 仿真 RTL，运行编译器生成的程序，结果与 L2 逐位比较。简报中标为“[M, sim]”的结果即来自这一级。L3 还提供每个节点的周期级时间线，用于校准性能模型。

\textbf{L4　硅片逐拍核对。} 在 VP815 上运行，读回 GDDR6 内容与 L2 逐位比较。按简报，完整 35 节点阵列的所有列出设计点都与编译器期望内存位真。

\subsection{单元级验证}

在整链验证之前，各硬件单元需要独立验证。本报告建议的单元级验证项如下 \tR：
\begin{itemize}
  \item \textbf{MLP72 模式}：INT8、SIGNED 4$\times$4、块浮点模式的乘加与累加行为，以随机向量与边界值（最大正负值、零、饱和）对照位真模型；
  \item \textbf{链}：不同 $W$、$d$、分块大小下的输出，以及行周期是否符合式~\eqref{eq:row}；
  \item \textbf{epilogue}：反量化、GELU/GeGLU、K 分段合并与向量节点同名运算逐位一致，覆盖查表的全部表项与插值边界；
  \item \textbf{向量节点}：每种运算的数值行为，特别是行归约运算在不同行长下的结果与溢出行为；
  \item \textbf{同步}：\texttt{WAIT}/\texttt{POST} 在各种到达顺序下的正确性，包括生产者先于消费者、消费者先于生产者、多个消费者等待同一标志；
  \item \textbf{跨时钟域}：以多组频率比运行异步 FIFO 与同步器的仿真，并做结构化的跨时钟域检查。
\end{itemize}

\subsection{系统级验证}

\textbf{整 chunk 回归。} 每一次 RTL 或编译器改动后，在 L3 与 L4 上运行完整 chunk，与 L2 逐位比较。由于整 chunk 仿真耗时较长，本报告建议分为三档 \tR：快速档只运行一层或一个积分步；标准档运行缩减的模型（例如减少层数）；完整档运行整 chunk。快速档用于每次提交，完整档用于每次比特流发布。

\textbf{多输入覆盖。} 46 帧录制数据覆盖了不同的视觉输入与机器人状态。由于设计没有数据依赖的执行路径（调度完全静态），多输入主要覆盖的是数值行为（例如激活幅值不同导致的饱和与尺度变化），而不是控制路径。

\textbf{故障定位。} 当 L4 与 L2 不一致时，逐拍核对可以定位到第一个出错的外存写入，进而定位到产生它的节点与指令。这使硅片调试的定位从“结果不对”缩小到“某节点某指令的输出不对”。

\subsection{回归基础设施}

逐位核对的价值在于它可以被自动化。本报告建议的回归基础设施包括 \tR：
\begin{itemize}
  \item \textbf{自动比对工具}：读入硅片或仿真的外存转储与期望内存，逐地址比较，输出第一个不一致的地址、所属张量、生产者节点与指令；
  \item \textbf{分档回归}：提交级（单层或单步）、夜间级（缩减模型）、发布级（整 chunk、全部 46 帧）三档；
  \item \textbf{结果数据库}：记录每次运行的构建版本、时延、精度与核对结果，形成随时间变化的性能与正确性曲线；
  \item \textbf{性能回归}：把每次运行的时延与上一次比较，超过阈值的退化自动报警。
\end{itemize}

\subsection{覆盖度量}

由于调度静态、命令流无数据依赖分支，传统的“代码覆盖率”在本设计中意义有限：整 chunk 回归一次就会执行所有被使用的命令路径。更有意义的覆盖度量是 \tR：
\begin{itemize}
  \item \textbf{模式覆盖}：每种链模式（INT8、INT4、块浮点、融合、分裂）与每种 epilogue 至少在一个整 chunk 中被使用并核对通过；
  \item \textbf{形状覆盖}：每种深度类、每种 $W$ 与分块组合至少出现一次，特别是边角分块（不满的最后一块）；
  \item \textbf{数值覆盖}：量化饱和、零尺度、查表边界等数值边界情况在单元测试中被覆盖；
  \item \textbf{同步覆盖}：不同的生产者—消费者到达顺序在仿真中被覆盖（可通过随机扰动节点启动时间实现）。
\end{itemize}

\subsection{硅片调试}

硅片上的问题可以分为三类：结果不一致、挂死（某节点永远等待）与性能异常。对应的调试手段是 \tR：
\begin{itemize}
  \item \textbf{结果不一致}：由自动比对工具定位第一个不一致的写入，再在 L3 仿真中复现该节点的执行；若仿真一致而硅片不一致，问题多在时序、跨时钟域或外存访问，而非逻辑；
  \item \textbf{挂死}：读取 GDDR6 中全部标志字的状态，找出正在等待的节点及其等待的标志，再回溯到本应写入该标志的生产者；
  \item \textbf{性能异常}：在各节点的关键事件（开始、每个 \texttt{WAIT} 的进入与退出、结束）处记录时间戳到外存，重建硅片上的时间线，与 L3 仿真的时间线和校准模型的预测对比。
\end{itemize}
其中时间戳记录同时服务于演示中的时延分解显示（第~\ref{sec:demo} 节）。

\subsection{性能模型的构建与校准}

简报中的时延构成来自一个“校准模型”\tP。从验证角度，这一模型的构建与校准可以分为三步：第一，按式~\eqref{eq:row} 与分块策略，从命令流推算每个节点每条命令的理想周期数；第二，引入外存访问时延、同步往返与跨时钟域开销等参数；第三，用 L3 仿真的周期级时间线与 L4 硅片的总时延校准这些参数。校准后的模型用于回答“如果……会怎样”的设计问题，例如运行时链分裂的收益（0.02\% 至 0.12\% \tP）与融合链对的时延收益（约 0 \tP）。模型的可信度需要持续以新的实测检验：每当一个由模型预测的设计点实际上板，就应比较预测与实测，并更新模型。

\subsection{性能验证}

性能验证比较校准模型的预测与 L3、L4 的实测。简报中的时延构成（向量约 50\%、链计算约 25\%、串行权重装载约 15\%、feeder 停顿约 6\%）就来自校准模型 \tP。模型的一个可检验预测是：运行时融合链对的读流量减少 11\% 至 23\% \tMsim{} 不会转化为时延收益（预计约 0 \tP \tB{B09}），因为当前瓶颈不在读带宽。融合链对的比特流构建完成后，其上板结果将直接检验这一预测。

\subsection{与校准模型的一致性}

校准模型在当前项目中承担两项职责：解释已有结果（给出时延构成），以及预测尚未实现的设计（例如链分裂与融合链对的收益）。其可信度可以从几组已有数据中得到侧面印证：GeGLU/GELU epilogue 在同一比特流上带来的 $-6.1\%$ \tM{} 与“向量工作占主导”的判断方向一致；K-split epilogue 在仿真中使对应向量阶段减少 81\%、整 chunk 减少 6.4\% \tMsim，也与“向量阶段在总时延中占比显著”的判断一致。更直接的检验将来自融合链对的上板：模型预测其时延收益约为 0 \tP，若实测也接近 0，将进一步确认模型对瓶颈的刻画。

\subsection{验证结果汇总}

表~\ref{tab:vresults} 汇总简报中的验证结果。

\begin{table}[t]
\centering
\caption{验证结果汇总（依据简报）。}
\label{tab:vresults}
\small
\begin{tabularx}{\textwidth}{l l Y}
\toprule
对象 & 验证级别 & 结果 \\
\midrule
首次整阵列运行 & L4 & 6.48~s，位真 \tM \\
INT8 最佳调度 & L4 & 1.221~s，位真；关节最大 $0.665^\circ$、平均 $0.072^\circ$ \tM \\
INT4 链 & L4 & 1.007~s；关节最大 $44.9^\circ$（不可用）\tM \\
INT4 + GeGLU/GELU epilogue & L4（同一比特流开关对比） & 1.047 对 1.115~s（$-6.1\%$），与 INT4 逐位相同 \tM \\
fabric-light 链 & L4 & 1.588~s，位真 \tM；时钟恢复后需重测 \tB{B13} \\
K-split 反量化 epilogue & L3 & 对应向量阶段 $-81\%$，整 chunk $-6.4\%$ \tMsim；比特流构建中 \tB{B08} \\
运行时融合链对 & L3 & 位真；读流量 $-11\%$ 至 $-23\%$ \tMsim；预计时延约 0 \tP；比特流构建中 \tB{B09} \\
链分裂 & L3 & 合法并已仿真；运行时切换比最佳固定形状好 0.02\% 至 0.12\% \tP \tB{B10} \\
BFP8 链 & L1 + L3 & 软件动作误差 0.70\%（INT8 为 0.97\%）\tMsw；仅 RTL 仿真 \tB{B11} \\
\bottomrule
\end{tabularx}
\end{table}

\subsection{验证环境与工具}

L3 级验证以 Verilator 把 RTL 编译为周期精确的 C++ 模型。为了在仿真中运行编译器生成的完整程序，仿真环境需要包括 \tR：
\begin{itemize}
  \item \textbf{外存模型}：以数组模拟 GDDR6 的内容，装入与硅片相同的初始镜像；其访问时延可以按固定值或简化的通道模型建模。外存模型越理想，越可能掩盖与外存时序相关的问题，这类问题需在 L4 上覆盖；外存模型的时延与带宽参数应以板上实测的 GDDR6 有效带宽校准 \tB{B04}。
  \item \textbf{NoC 模型}：以简化的请求—响应模型代替真实 NoC，按 NAP 接口的协议与节点交互。
  \item \textbf{主机模型}：以脚本完成输入写入、启动标志与完成检测，与真实主机程序使用同一套接口定义。
  \item \textbf{比对与时间线}：仿真结束后转储外存，与期望内存比对；同时记录每个节点的事件时间线，用于性能校准。
\end{itemize}
由于整 chunk 包含大量周期，完整的 L3 仿真耗时较长；分档回归中的缩减模型（减少层数或积分步数）可以在保持覆盖的前提下显著缩短仿真时间。

\subsection{错误注入测试}

为确认验证体系本身能够发现错误，本报告建议定期进行错误注入测试 \tR：在编译器输出中故意删除一对 \texttt{WAIT}/\texttt{POST}、改变一个量化尺度、或交换两个分块的地址，确认 L3 或 L4 的比对能够报告不一致，并且自动比对工具能把不一致定位到被修改的位置。这类测试验证的是“验证链的灵敏度”，防止出现比对工具本身失效而无人察觉的情况。

\subsection{典型问题类型与对应的验证级别}

根据验证链的结构，可以预先判断不同类型的问题会在哪一级被发现，从而把验证资源集中在最有效的环节 \tR：
\begin{itemize}
  \item \textbf{数值方案问题}（例如某层量化尺度选择不当导致误差过大）：在 L0 对 L1 的精度比较中发现，与硬件无关；
  \item \textbf{调度改变数值}（例如 K 分段合并顺序与位真模型不一致）：在 L1 对 L2 的比较中发现；
  \item \textbf{RTL 逻辑错误}（例如 epilogue 的舍入与向量节点不一致）：在 L2 对 L3 的比较中发现，通常可以在单元级测试中更早暴露；
  \item \textbf{同步遗漏}：在 L3 中可能因时序巧合而不暴露，在 L4 上以偶发的不一致或挂死表现；随机扰动节点启动时间的 L3 仿真有助于提前发现；
  \item \textbf{时序与跨时钟域问题}：L3 仿真通常不能发现，只在 L4 上以不可复现的错误表现，需要依靠时序签核与跨时钟域检查预防；
  \item \textbf{外存与 NoC 相关问题}：L3 仿真中的外存模型若过于理想，可能掩盖这类问题，需要在 L4 上以压力测试覆盖。
\end{itemize}
这一分析说明，逐拍核对虽然能可靠地\emph{发现}问题，但对“L3 通过、L4 失败”的问题，定位仍需依靠时序与跨时钟域方面的专门手段。

\subsection{精度验证的统计考虑}

46 帧留出数据上的关节最大误差 $0.665^\circ$ 与平均误差 $0.072^\circ$ \tM{} 是当前的精度结论。从统计上看，需要注意以下几点 \tR：
\begin{itemize}
  \item \textbf{样本量}：46 帧对于判断“平均行为”已有一定代表性，但对于“最坏情况”的估计有限；最大误差可能随样本增加而上升。建议在扩大样本（例如更多留出帧）后报告最大误差的变化。
  \item \textbf{误差分布}：除最大值与平均值外，建议报告每帧最大误差的分布（例如分位数），以及误差在动作块内随时间步的变化，判断误差是否随积分或时间步累积。
  \item \textbf{按关节分解}：7 个关节的误差可能差异较大，按关节报告有助于判断误差对末端位置的影响。
  \item \textbf{与参考的差异来源}：若个别帧的误差显著偏大，应检查这些帧的激活是否出现饱和或尺度异常，这可以借助位真模型逐层定位。
\end{itemize}
由于硅片与位真模型逐位一致，上述统计分析都可以在软件中完成，不需要额外的硅片运行。

\subsection{各设计点的验证状态}

综合表~\ref{tab:vresults}，各设计点的验证状态可以归为三类：
\begin{itemize}
  \item \textbf{已完成硅片验证且精度可用}：INT8 最佳调度；fabric-light 链（位真，精度与 INT8 相同的数值方案）。
  \item \textbf{已完成硅片验证但精度不可用}：INT4 链及其 epilogue 变体。这些设计点的价值在于验证了 INT4 模式与 epilogue 的硬件正确性，以及 epilogue 的结构收益（$-6.1\%$ \tM），这一收益有望迁移到 INT8 设计点。
  \item \textbf{仅完成仿真}：K-split epilogue、融合链对、链分裂、BFP8 链。其中 K-split 与融合链对的比特流正在构建。
\end{itemize}
需要特别注意，epilogue 的 $-6.1\%$ 是在 INT4 比特流上测得的；它在 INT8 设计点上的收益需要在相应的比特流上重新测量，不能直接套用 \tB{B12}。

\subsection{验证策略的经验}

从当前的验证结果可以归纳出几条对类似项目有参考价值的经验：
\begin{itemize}
  \item \textbf{先建位真模型，再写 RTL}：位真模型既是精度评估的工具，也是期望内存的来源；它越早确定，RTL 与编译器的开发就越有明确的正确性目标。
  \item \textbf{让编译器输出期望值}：把期望内存作为编译器的标准输出，使每一次调度改动都自带验证数据，这是调度能被积极优化（1.901~s $\to$ 1.221~s \tM）的前提。
  \item \textbf{同一比特流上做开关对比}：epilogue 的收益（$-6.1\%$ \tM）之所以可信，是因为它是在同一比特流、同一时钟下开关得到的；跨比特流的比较混入了时钟与布局的影响。
  \item \textbf{负面结果同样有价值}：INT4 设计点在硬件上是正确的（与期望内存一致），但精度不可用；这一结论只有在“正确性”与“精度”被干净分离的验证体系下才能得出。
\end{itemize}

\subsection{验证缺口}

对照验证目标，目前存在以下缺口，本报告建议逐项补齐 \tR：
\begin{enumerate}
  \item \textbf{闭环精度}：46 帧开环误差不能代表闭环成功率，需要在仿真或真实机器人上补充闭环评测。
  \item \textbf{可用性阈值}：“可用”与“不可用”之间尚无量化阈值，建议以任务需求为据给出关节误差的验收上限。
  \item \textbf{时延分布}：简报给出的是单次或典型时延，缺少多次运行的时延分布（方差、高分位数），无法以测量支撑“确定性时延” \tB{B02}。
  \item \textbf{能耗}：缺少板卡功耗与每 chunk 能耗的测量 \tB{B03}。
  \item \textbf{Thor 端对照细节}：缺少 Thor 端相对 fp32 的误差与时延分解，使“匹配精度”的比较口径不完全对等 \tB{B14}。
  \item \textbf{主机开销}：简报给出约 47~ms/次调用，未注明实测或预测，需要确认并按环节分解 \tB{B05}。
  \item \textbf{长时间稳定性}：缺少长时间连续运行的挂死、核对失败与时延漂移统计 \tB{B15}。
  \item \textbf{同步开销}：“feeder 停顿约 6\%”\tP{} 缺少实测分解 \tB{B16}。
  \item \textbf{分阶段时延}：缺少视觉编码、前缀与动作专家单步的分阶段实测时延，限制了对减步数、token 复用等方向的评估 \tB{B01}。
\end{enumerate}


%======================================================================
\section{待上板测量清单}\label{sec:board}
%======================================================================

本节汇总本报告中\textbf{尚需在硬件上实测}的数据。正文中以 \tB{Bxx} 在验证、性能、演示与后续计划章节中标出这些数据出现的位置；设计方案章节不作标记。清单分为三类：\emph{测量补齐}（现有比特流即可测，不需要改动硬件）；\emph{新比特流上板}（需要尚在构建或尚待设计的比特流）；\emph{对照平台}（在 Jetson AGX Thor 上测量）。与之相对，凡是“硅片与位真模型逐位一致”即可推出的精度结论（例如扩大样本后的误差统计、逐层敏感度分析），都可以在软件中完成，不列入本清单。

\begin{small}
\begin{longtable}{p{0.8cm} p{2.6cm} p{4.3cm} p{3.3cm} p{2.4cm}}
\caption{待上板测量清单。“现状”一栏为简报中已有的数据或其来源类型；“判据”为本报告建议的通过条件或用途 \tR。}\label{tab:board}\\
\toprule
编号 & 测量项 & 方法与条件 & 现状 & 判据 / 用途 \\
\midrule
\endfirsthead
\toprule
编号 & 测量项 & 方法与条件 & 现状 & 判据 / 用途 \\
\midrule
\endhead
\bottomrule
\endfoot
\multicolumn{5}{l}{\textit{A. 测量补齐：现有最佳 INT8 比特流（500 / 288.9 / 216.7~MHz，35 节点）即可测}} \\
B01 & 分阶段时延 & 在各节点关键事件处写时间戳（启动、每个 WAIT 进出、结束），重建 SigLIP、前缀、动作专家每一步的时间线 & 无实测，仅有整 chunk 1.221~s \tM & 评估减步数、token 复用、KV 驻留的收益 \\
B02 & 时延分布 & 同一输入连续运行不少于数百次，另以 46 帧各运行多次；分别记录阵列时延与端到端时延 & 无 & 给出均值、标准差、P99、最大值；支撑“确定性时延” \\
B03 & 功耗与每 chunk 能耗 & 板卡供电口测量空闲、运行时平均与峰值功耗；乘以时延得每 chunk 能耗；长时间运行记录芯片温度 & 无 & 与 Thor 做能效对比；确认散热条件下时钟可维持 \\
B04 & GDDR6 有效带宽 & 以微基准测试单节点与多节点并发经 NoC 的读写带宽与时延，覆盖条带化与非条带化布局 & 无 & 校准性能模型中的 $BW_{\text{eff}}$ 与外存时延参数 \\
B05 & 主机开销 & 分别计时输入写入、启动、完成检测与读回；区分 PCIe 传输与主机软件开销 & 约 47~ms/次调用（简报未注明实测或预测） & 确认数值；为百毫秒目标定位主机侧优化点 \\
B06 & 时延构成校核 & 由 B01 的时间戳统计向量、链计算、串行权重装载、feeder 停顿的时间占比 & 50 / 25 / 15 / 6\% \tP & 与校准模型偏差应在可解释范围内，否则更新模型 \\
B07 & INT4 设计点关键路径构成 & 在 INT4 比特流上重复 B01、B06 & 向量约 82\% 关键路径 \tP & 检验“GEMM 加速后向量主导”的判断 \\
B15 & 长时间稳定性 & 连续运行数小时的 46 帧循环，统计挂死、核对失败与时延漂移 & 无 & 零挂死、零核对失败；验收标准第 5 条 \\
B16 & 同步与 feeder 停顿 & 由时间戳统计每次 WAIT 的等待时间及其分布，区分同步往返与数据未就绪 & 约 6\% \tP & 判断标志字放置与重叠策略的改进空间 \\
B17 & 演示口径时延 & 在演示主机上测阵列时延与端到端时延 & 无 & 演示界面显示的数据来源 \\
\multicolumn{5}{l}{\textit{B. 新比特流上板：需要构建中或待设计的比特流}} \\
B08 & K-split 反量化 epilogue & 同一比特流上开 / 关该 epilogue，测整 chunk 时延与对应向量阶段时间；核对位真；测时钟 & 向量阶段 $-81\%$、整 chunk $-6.4\%$ \tMsim；构建中 & 位真；实测收益与仿真偏差可解释 \\
B09 & 运行时融合链对 & 同一比特流上开 / 关融合，测读流量（经时间戳或 NoC 计数）与整 chunk 时延 & 读流量 $-11\%$ 至 $-23\%$ \tMsim；时延约 0 \tP；构建中 & 检验校准模型的“约 0 收益”预测 \\
B10 & 链分裂（可选） & 同一比特流上比较运行时切换与固定形状的时延 & 0.02\% 至 0.12\% \tP；仅仿真 & 仅在需要检验模型时进行 \\
B11 & BFP8 链 & 构建 BFP8 模式比特流；测 46 帧精度（与位真模型一致性）、时钟、资源与整 chunk 时延 & 软件误差 0.70\% 对 INT8 0.97\% \tMsw；仅 RTL 仿真 & 位真；精度不劣于 INT8；记录时延变化 \\
B12 & INT8 上的 GeGLU/GELU epilogue & 在 INT8 比特流上加入该 epilogue，做同一比特流开 / 关对比 & 仅在 INT4 比特流上测得 $-6.1\%$ \tM & 确认收益能否迁移到 INT8 设计点 \\
B13 & fabric-light 时钟恢复 & 取得该设计点时序报告；实施分段寄存等措施后重新上板，测三域时钟与整 chunk 时延 & 310 / 216.7 / 177.3~MHz，1.588~s \tM & 判断能否作为后续基线结构（第~\ref{sec:risk} 节决策点 1） \\
\multicolumn{5}{l}{\textit{C. 对照平台：Jetson AGX Thor}} \\
B14 & Thor 对照细节 & 以同一 46 帧运行 bf16 版本；测相对 fp32 的关节误差、时延分布、分阶段时延与功耗模式 & 仅有 125.3~ms \tM & 使“匹配精度”的对比口径对等 \\
\end{longtable}
\end{small}

\textbf{优先级建议 \tR。} 第一批为 B01、B02、B03、B05、B06：它们只需要现有比特流与测量脚本，却决定了后续所有设计决策的依据。第二批为 B08、B09（对应已在构建中的比特流，其结果直接检验校准模型）与 B12（需在 INT8 比特流上加入已验证过的 epilogue）。第三批为 B11、B13 与 B14。B04、B15、B16、B17 可以与第一批的测量脚本共用基础设施，穿插进行。

\textbf{测量的共同要求 \tR。} 所有上板测量都应记录完整的构建产物版本（比特流、硬件描述、编译器、程序镜像、期望内存）与测量条件（时钟、主机配置、输入帧编号）；每次测量都应同时完成逐位核对，确保性能数据来自结果正确的运行。

%======================================================================
\section{性能分析}\label{sec:perf}
%======================================================================

\subsection{时延构成}

\begin{figure}[t]
\centering
\begin{tikzpicture}
\begin{axis}[
  xbar, width=0.7\textwidth, height=4.6cm,
  xmin=0, xmax=60, xlabel={占 1.221~s 的比例（\%）},
  symbolic y coords={feeder 停顿, 串行权重装载, 链计算, 向量工作},
  ytick=data, nodes near coords, nodes near coords align={horizontal},
  bar width=9pt, enlarge y limits=0.2, font=\small]
\addplot[fill=blue!35] coordinates {(6,feeder 停顿) (15,串行权重装载) (25,链计算) (50,向量工作)};
\end{axis}
\end{tikzpicture}
\caption{最佳 INT8 设计点（1.221~s）的时延构成 \tP，依据简报的校准模型。}
\label{fig:breakdown}
\end{figure}

图~\ref{fig:breakdown} 给出最佳 INT8 设计点的时延构成 \tP，需以板上时间戳实测校核 \tB{B06}。换算为时间 \tD：向量工作约 0.61~s，链计算约 0.31~s，串行权重装载约 0.18~s，feeder 停顿约 0.07~s，其余约 0.05~s 未在简报中归类。

\subsection{阵列利用率}

等效 MAC 率为 $1.47~\text{TMAC}/1.221~\text{s}\approx1.20$~TMAC/s \tD。若按 2,560 个 MLP72 $\times$ 16 INT8 MAC/周期 $\times$ 500~MHz 计，全器件峰值约 20.5~TMAC/s，等效利用率约 6\% \tD；若只计实际使用的约 760 至 860 个 MLP72，其峰值上界约 6.1 至 6.9~TMAC/s，等效利用率约 17\% 至 20\% \tD（其中包含 feeder 等非乘法级，实际乘法级利用率会更高）。作为对照，Jetson AGX Thor 的 125.3~ms 对应约 11.7~TMAC/s 的等效率 \tD。

这组数字给出的结论是：与 Thor 的差距主要不来自“乘法器不够”，而来自“乘法器之外的时间太多”（向量与权重装载约 65\% \tP）和“能布下的乘法器太少”（约三分之一）。

\subsection{设计演进}

在同一套 INT8 硬件上，调度把时延从 1.901~s 降到 1.221~s；硬件结构上沿三条线探索：降低位宽（INT4，快约 18\% 但精度不可用 \tD）、把向量工作移入链（epilogue，每项约 6\%）、缓解布线约束（fabric-light，布下更多乘法器但时钟下降）。运行时可重构（融合链对、链分裂）的预期收益接近零 \tP。这些结果与时延构成的判断一致：在向量与权重装载成为主导之前，GEMM 侧的优化收益会被 Amdahl 定律限制。

\subsection{性能敏感性}

以时延构成 \tP{} 为基础，可以计算各项对总时延的边际敏感性 \tE：若某一项时间减少 10\%，总时延相应减少该项占比的 10\%。据此：
\begin{itemize}
  \item 向量工作每减少 10\%，总时延约减少 5\%（约 61~ms）；
  \item 链计算每减少 10\%，总时延约减少 2.5\%（约 31~ms）；
  \item 串行权重装载每减少 10\%，总时延约减少 1.5\%（约 18~ms）；
  \item feeder 停顿每减少 10\%，总时延约减少 0.6\%（约 7~ms）。
\end{itemize}
这一简单的敏感性排序与第~\ref{sec:risk} 节的路线排序一致：同样的工程投入，花在向量通路上的收益约为花在链计算上的两倍。需要注意，这一计算假设各项独立、重叠关系不变；当某一项被大幅压缩后，其他项可能暴露出来，敏感性排序会随之改变，因此每一轮优化之后都需要重新测量时延构成。

\subsection{时延构成随设计点的变化}

简报给出了两组时延构成信息：最佳 INT8 点的时间构成（向量约 50\%）与 INT4 点的关键路径构成（向量约 82\%）\tP，均需以板上时间戳实测校核 \tB{B06}\tB{B07}。尽管口径不同，二者揭示的趋势是一致的：当链计算被加速（INT4 使每级乘法吞吐翻倍），向量工作在关键路径中的比重显著上升。这正是 Amdahl 定律在本设计中的表现，也解释了为什么 INT4 的时延只比 INT8 低约 18\% \tD：乘法吞吐翻倍，但总时延中只有约四分之一属于链计算。可以预期，随着向量通路优化的推进，链计算与权重装载的比重将重新上升，届时 BFP8、fabric-light 与融合链对等 GEMM 侧优化的价值才会显现。

\subsection{情景推算}

以下情景基于简单的加性模型：假设各项时间相互独立、优化某一项不改变其他项、主机开销不变。该模型偏乐观，结果全部标为 \tE，只用于比较方向：
\begin{itemize}
  \item 仅向量通路加速 2 倍：约 0.91~s；
  \item 向量加速 4 倍，且权重装载完全隐藏：约 0.58~s；
  \item 在上一条基础上链计算再加速 2 倍：约 0.43~s，仍约为 Thor 参照的 3 倍以上。
\end{itemize}
这说明向量通路与权重装载是必须首先解决的问题；要在单片器件上追平 Thor，还需要算法侧配合（例如减少积分步数或视觉 token 数）或多器件扩展。

\subsection{与 GPU 平台的对比讨论}

与 Jetson AGX Thor 的对比需要放在正确的框架中理解。Thor 端的 125.3~ms \tM{} 是以 bf16 编译的结果，依托成熟的 GPU 编译栈，自动完成了算子融合与存储管理；本设计则从 RTL 到编译器完全自研，对每个算子的数据流都有完全的控制。二者的差距约 9.7 倍 \tD，可以从三个层面理解。

第一，\emph{有效算力}：FPGA 端的等效 MAC 率约 1.20~TMAC/s，Thor 端约 11.7~TMAC/s \tD，但这一对比混合了乘法与非乘法时间。若只看乘法时间（链计算约 0.31~s \tD），FPGA 端在链计算期间的等效率约 4.7~TMAC/s \tD（$1.47/0.31$），差距缩小到约 2.5 倍。这说明乘法阵列本身并不是主要差距来源。

第二，\emph{非乘法开销}：FPGA 端约 65\% 的时间 \tP{} 花在向量运算与权重装载上。GPU 通过 kernel 融合与大容量缓存把大部分非线性运算藏在矩阵乘的访存阴影中，而本设计目前以独立的向量遍历执行它们。

第三，\emph{数值格式}：Thor 端为 bf16，本设计为 W8A8。INT8 本应带来吞吐优势，但它引入了每个 GEMM 前后的量化与反量化，而这部分工作恰好落在向量通路上。这是一个值得重视的反直觉结论：在向量通路受限的设计中，低位宽的收益会被量化开销部分抵消。

需要强调，Thor 端的时延分解与精度数据简报均未给出，上述讨论只能基于 FPGA 端的数据与一般性的 GPU 知识，不能据此推断 Thor 端的具体构成。

\subsection{对设计决策的反思}

回顾当前的设计与结果，可以对几项关键决策做出评价。

\textbf{列并行链：正确但代价在布线。} 它使每一级的乘法器满载，这一目标已经达到；但它把大量分布式连线引入 fabric，最终使布线先于乘法器耗尽。fabric-light 链证明这一代价可以大幅降低，下一步的关键是恢复其时钟。

\textbf{矩阵与向量分离：简化了设计，但低估了向量工作。} 分离使链保持简单、使向量运算可以用浮点精确实现，并使位真模型清晰。但 16 条向量 lane 与 24 条 GEMM 链的配比明显偏向矩阵侧。epilogue 是对这一失衡的修补；更根本的是重新配比。

\textbf{外存为枢纽：换来了可扩展性与可验证性。} 节点间只经外存通信，使同步语义简单、使期望内存可以精确计算，也使节点数可以按需增减。它的代价（同步往返与中间结果往返）目前约占 feeder 停顿的 6\% \tP{} 以及向量工作中的一部分往返，尚在可接受范围内。

\textbf{编译器中心：最有价值的投入。} 调度带来的 1.56 倍改善 \tD{} 与逐拍核对能力，都是编译器提供的。对于一个需要快速迭代硬件设计点、又必须保证结果正确的项目，这是最有杠杆的投入。

\subsection{与目标的差距分解}

以 Thor 参照（125.3~ms \tM）为目标，当前最佳可用点（1.221~s \tM）需要缩短约 1.10~s \tD。按时延构成 \tP{} 分解，这一差距中：向量工作约 0.61~s，链计算约 0.31~s，串行权重装载约 0.18~s，feeder 停顿约 0.07~s，未归类约 0.05~s \tD。即使向量工作与权重装载全部消除，剩余的链计算与停顿仍约 0.38~s，是参照的约 3 倍 \tD。这再次说明，追平参照需要多个方向同时推进：向量通路、权重驻留、更高的阵列利用率与时钟，以及算法侧的配合。这一分解也为各方向设定了量化目标：例如，若目标是 0.8~s，则在其他各项（合计约 0.61~s）不变时，向量工作需要降到约 0.19~s 以下 \tE；若目标是 0.5~s，则仅优化向量工作已不足以达到，必须同时压缩链计算、权重装载与未归类开销。

%======================================================================
\section{演示方案}\label{sec:demo}
%======================================================================

\subsection{演示目标}

演示的目的是让评审者在现场直观看到三件事：\emph{完整的 \pizero{} 推理确实全部在单片 FPGA 上完成}；\emph{FPGA 输出的动作与 fp32 参考高度一致}；\emph{每一次运行的结果都可以被逐拍核对}。演示以最佳 INT8 设计点（完整 35 节点阵列，500 / 288.9 / 216.7~MHz）为基础，不使用精度不可用的 INT4 设计点。

\subsection{系统组成}

\begin{figure}[t]
\centering
\begin{tikzpicture}[font=\footnotesize, >=Stealth, node distance=8mm,
  bx/.style={draw, rounded corners=2pt, fill=#1, text width=2.8cm, align=center, minimum height=12mm},
  bx/.default=white]
\node[bx=gray!10] (data) {DROID 录制帧\\（2 路图像 + 指令 + 状态）};
\node[bx=blue!8, right=of data] (host) {主机\\输入装载、启动、读回、核对};
\node[bx=green!10, right=of host] (fpga) {VP815 / AC7t1500\\35 节点阵列\\完整 \pizero{} chunk};
\node[bx=yellow!15, below=of host] (ui) {显示界面\\动作轨迹对比、误差、时延、核对状态};
\node[bx=gray!10, below=of data] (ref) {fp32 参考结果\\（预先计算）};
\draw[->] (data) -- (host); \draw[<->] (host) -- node[above]{PCIe} (fpga);
\draw[->] (host) -- (ui); \draw[->] (ref) -- (ui);
\end{tikzpicture}
\caption{演示系统组成。}
\label{fig:demo}
\end{figure}

演示系统由四部分组成（图~\ref{fig:demo}）：
\begin{enumerate}
  \item \textbf{数据源}：46 帧 DROID 留出录制数据，每帧包含两路相机图像、语言指令与机器人状态；
  \item \textbf{主机}：负责把输入写入 GDDR6、启动阵列、等待完成、读回动作块与校验区域；
  \item \textbf{FPGA 板卡}：VP815 上的 AC7t1500，运行完整 35 节点阵列，执行完整 \pizero{} chunk；
  \item \textbf{显示界面}：展示输入图像与指令、FPGA 输出的 50 步动作轨迹与 fp32 参考轨迹的对比、逐关节误差、每次推理的时延、以及逐位核对状态。
\end{enumerate}

\subsection{演示环境要求}

演示需要以下环境 \tR：一台安装 VP815 板卡的主机（PCIe Gen4 x16 插槽，满足板卡的供电与散热要求）；主机上的驱动、主机程序与显示界面；预先准备的构建产物与 46 帧数据；一台用于投影的显示器。若采用扩展演示中的运动学可视化，还需要相应的仿真软件。所有软件与数据应在演示前固定版本并完成一次完整彩排。

\subsection{演示数据的选择}

演示优先使用 46 帧留出数据中的代表帧 \tR：选择场景与指令差异较大的若干帧，使评审者看到模型在不同输入下的输出；同时准备误差最大的那一帧，主动展示最坏情况下 FPGA 输出与参考输出的差异，以增强结果的可信度。所有帧的选择与顺序在演示前固定，避免现场临时挑选带来的偏差。

\subsection{演示流程}

\begin{enumerate}
  \item \textbf{初始化}：把程序、权重与常数表装入 GDDR6；显示所用比特流的设计点信息（时钟、节点数、数值方案）。
  \item \textbf{单帧推理}：选择一帧录制数据，显示其图像与指令；启动推理；完成后显示动作轨迹、与 fp32 参考的逐关节误差，以及端到端时延（含主机开销）。
  \item \textbf{逐位核对}：读回校验区域，与编译器期望内存比较，显示“位真”状态；若评审者希望，可以现场查看任意中间张量的核对结果。
  \item \textbf{批量回放}：依次运行全部 46 帧，汇总关节最大误差与平均误差，与报告中的 $0.665^\circ$ 与 $0.072^\circ$ \tM{} 对照；同时统计每帧时延，给出时延分布。
  \item \textbf{对照说明}：展示 Jetson AGX Thor 上同一 chunk 的 125.3~ms \tM{} 参照，以及本设计的时延构成（图~\ref{fig:breakdown}），说明差距来源与后续路线。
\end{enumerate}

第 4 步的批量回放同时补齐了第~\ref{sec:verif} 节列出的“时延分布”缺口 \tB{B02}：多次运行的时延统计可以直接作为确定性的实测证据。

\subsection{演示准备清单}

为保证演示的可重复性，演示前应完成以下准备 \tR：
\begin{itemize}
  \item 固定一组构建产物（比特流、硬件描述、编译器版本、程序镜像、期望内存），并在演示机上完成至少一次全部 46 帧的回放与核对；
  \item 预先计算并保存 46 帧的 fp32 参考动作与期望内存，演示时不在现场重新生成；
  \item 准备时延的两种口径：阵列时延（从启动标志到完成标志）与端到端时延（含主机装载与读回），界面上同时显示 \tB{B17}；
  \item 准备一份一页纸的说明，列出设计点参数、精度结果、时延与 Thor 参照，所有数字标注来源（实测或模型预测）。
\end{itemize}

\subsection{界面设计要点}

演示界面建议分为四个区域 \tR：左上为当前帧的两路相机图像与语言指令；右上为 50 步动作块中各关节角随时间的曲线，FPGA 结果与 fp32 参考叠加显示；左下为逐关节的最大与平均误差，以及 46 帧的累计统计；右下为时延信息与核对状态，包括本次阵列时延、端到端时延、历次时延的直方图，以及“位真 / 不一致”的状态指示。界面上所有数字的含义与来源都应清楚标注，避免把模型预测的时延构成与实测的总时延混为一谈。

\subsection{可选的扩展演示}

在基础演示之外，本报告建议两项可选扩展 \tR：
\begin{itemize}
  \item \textbf{动作可视化}：把 FPGA 输出的关节角序列送入机械臂运动学仿真，以三维动画并排显示 FPGA 与 fp32 参考的轨迹，使 $0.665^\circ$ 量级的误差在视觉上可感知。
  \item \textbf{同步开关对比}：在带 epilogue 的比特流上，现场对比开启与关闭 epilogue 的时延，展示“同一比特流、逐位相同、时延降低”的效果。
\end{itemize}
闭环真实机器人演示需要以秒级时延支撑实时控制，在当前时延水平下不作为本阶段的演示内容；待时延进入百毫秒量级后再规划。

\subsection{讲解要点}

演示时建议围绕以下三个要点讲解，每个要点都有界面上的直接证据 \tR：
\begin{enumerate}
  \item \textbf{“全部在 FPGA 上”}：展示主机侧程序只有装载、启动、等待与读回四个步骤；展示设计点信息中的 35 个节点与其分工。
  \item \textbf{“结果可信”}：展示 FPGA 动作曲线与 fp32 参考曲线几乎重合，逐关节误差在 1 度以内；展示逐位核对的“位真”状态。
  \item \textbf{“瓶颈清楚、路线明确”}：展示时延构成中向量工作约占一半 \tP，说明后续优化的方向；展示与 Thor 参照的差距，并说明情景推算给出的改进空间。
\end{enumerate}
讲解中应明确区分实测数据与模型预测：总时延与精度为实测，时延构成为校准模型预测。

\subsection{演示后的延伸评测}

演示本身会产生一组有价值的数据：46 帧的时延分布与逐帧精度。建议在演示后把这组数据整理为正式的评测记录 \tR，并在此基础上补充两项评测：一是在同一台主机上测量板卡功耗，计算每 chunk 能耗 \tB{B03}；二是以同一组 46 帧在 Thor 上运行 bf16 版本 \tB{B14}，记录其相对 fp32 的误差与时延分布，使两个平台的对比在精度与稳定性两个维度上都完整。

\subsection{评审中可能的问题与回答要点}

为使演示后的讨论更有效，本节列出几个可能被问到的问题及其回答要点，回答均以简报事实为据：
\begin{itemize}
  \item \textbf{问：为什么比 Thor 慢这么多？} 答：按校准模型，约 65\% 的时间 \tP{} 花在向量运算与权重装载上，乘法阵列本身不是主要差距；而且布线限制使只有约三分之一的 MLP72 在使用。后续路线以向量通路与片上驻留为先。
  \item \textbf{问：为什么不用 INT4？} 答：INT4 设计点在硬件上是正确的，时延 1.007~s，但关节最大误差 $44.9^\circ$ \tM，不可用；后续以 BFP8 与逐层敏感度分析为方向。
  \item \textbf{问：“位真”到底意味着什么？} 答：硅片上每一个写入外存的数据都与编译器按位真模型计算的期望值逐位相同；因此硅片的精度等于位真模型的精度，任何不一致都能定位到具体节点与指令。
  \item \textbf{问：为什么不做可重构阵列？} 答：已实现运行时融合链对与链分裂；按校准模型，前者时延收益约为 0、后者只有 0.02\% 至 0.12\% \tP，因为当前瓶颈不在阵列形状适配。
  \item \textbf{问：这个架构能换别的模型吗？} 答：结构相近的模型主要是重新编译；引入新算子时需要扩展向量节点与位真模型。
\end{itemize}

\subsection{演示风险与预案}

\begin{itemize}
  \item \textbf{比特流或板卡故障}：准备已验证的备用比特流与程序镜像；准备同一流程在 Verilator 仿真上的录屏作为后备。
  \item \textbf{主机开销波动}：显示时延时分别给出阵列时延与端到端时延，避免主机侧波动掩盖阵列的确定性。
  \item \textbf{核对失败}：一旦出现核对失败，界面应直接给出第一个不一致的外存地址与对应节点，这本身也是验证体系能力的展示。
\end{itemize}

%======================================================================
\section{风险与后续计划}\label{sec:risk}
%======================================================================

\subsection{风险评估}

\begin{table}[h]
\centering
\caption{主要风险。}
\label{tab:risk}
\small
\begin{tabularx}{\textwidth}{l Y Y}
\toprule
风险 & 表现 & 应对 \\
\midrule
向量通路瓶颈 & 约 50\% 时延 \tP；INT4 后约 82\% 关键路径 \tP & 结构性扩容、提频、epilogue 扩展 \\
布线约束 & 约三分之一 MLP72 可用；更大阵列无法布线 & fabric-light 推广；区域约束与流水化 \\
精度风险 & INT4 不可用 \tM；BFP8 仅软件数据 \tMsw & 逐层敏感度分析；BFP8 上板验证 \\
时钟下降 & fabric-light 阵列域 310~MHz & 时序收敛措施；与向量扩容联动评估 \\
收益为零的投入 & 融合链对、链分裂预期收益接近零 \tP & 定位为瓶颈转移后的储备 \\
可比性 & Thor 端缺少精度与分解数据 & 补充对照测量 \\
评估面 & 仅有开环误差 & 补充闭环评测 \\
\bottomrule
\end{tabularx}
\end{table}

\subsection{外部依赖风险}

项目还依赖若干外部条件，其变化可能影响进度 \tR：工具链版本（ACE 10.5.2 与 Synplify）的升级可能改变布局布线结果与可达时钟，升级前应以现有设计点做回归；板卡与驱动的稳定性直接影响硅片验证与演示；模型检查点的更新（例如换用新的微调检查点）需要重新生成量化参数、期望内存与精度基线。建议把这些外部依赖的版本也纳入构建产物的可追溯记录。

\subsection{后续设计路线}

按预期收益排序，后续设计路线如下 \tR：
\begin{enumerate}
  \item \textbf{向量通路扩容与提频}：若向量工作整体加速 2 倍，上界约节省 0.31~s \tD（$1.221\times50\%\times1/2$）。代价是与 GEMM 链争夺布线，需要在校准模型中扫描“链数 $\times$ 向量 lane 数 $\times$ 时钟”的组合后再综合。
  \item \textbf{epilogue 扩展}：在已上板的 GeGLU/GELU 与已仿真的 K-split 之外，把反量化与残差、按 token 量化、softmax 的部分运算移入链或 PV 链的结果通路。
  \item \textbf{BFP8 上板}：以精度为主要收益，同时评估对尺度处理的简化。
  \item \textbf{片上驻留}：前缀 KV（INT8 约 4.8~MB \tE）与部分动作专家权重驻留片上，削减串行权重装载（上界约 0.18~s \tD）。
  \item \textbf{主机开销优化}：减少每次调用的主机参与，为百毫秒量级目标做准备。
  \item \textbf{测量补齐}：分阶段时延、能耗、时延分布、Thor 端对照细节与闭环评测（见第~\ref{sec:board} 节清单 B01--B06、B14）。
\end{enumerate}

\subsection{已知限制}

在进入后续阶段之前，需要明确当前设计的已知限制：
\begin{itemize}
  \item 最佳可用设计点的时延（1.221~s）距离闭环实时控制的需求仍有一个数量级以上的差距；
  \item 低于 INT8 的均匀量化不可用，低位宽路线需要以逐层敏感度分析为前提；
  \item 阵列规模受布线限制，约三分之一的 MLP72 在使用；
  \item 缺少功耗、时延分布与闭环精度数据；
  \item 当前设计与 Speedster7t 的 MLP72 级联、二维 NoC 和 GDDR6 深度绑定，移植到其他器件需要重做数据流与编译器后端。
\end{itemize}

\subsection{决策点}

后续推进中有几个需要明确做出的决策，本报告给出建议 \tR：
\begin{enumerate}
  \item \textbf{基线结构是否转向 fabric-light？} 建议以“在 fabric-light 上恢复时钟”为前提条件：若阵列与向量时钟能恢复到接近最佳 INT8 点的水平，则转向 fabric-light，以其更多的乘法器与更低的布线占用为后续扩容留出空间；否则保持原结构。
  \item \textbf{继续投入运行时可重构吗？} 建议暂停新的投入，只完成已在构建中的融合链对比特流的上板测量，用于检验校准模型；待向量通路优化使瓶颈转移后再复评。
  \item \textbf{低位宽路线如何走？} 建议放弃均匀 INT4，转向“BFP8 上板 + 逐层敏感度分析”，以精度为首要目标，速度收益作为次要目标。
  \item \textbf{何时做闭环机器人演示？} 建议以时延进入百毫秒量级、并完成仿真闭环评测为前提。
\end{enumerate}

\subsection{里程碑}

后续工作建议按以下里程碑推进 \tR：
\begin{enumerate}
  \item \textbf{M1　测量补齐}：分阶段时延、时延分布、板卡功耗与每 chunk 能耗、Thor 端对照细节（清单 B01--B07、B14--B17）。这是其余决策的依据，且不需要改动硬件。
  \item \textbf{M2　已在构建中的比特流上板}：K-split 反量化 epilogue 与运行时融合链对的比特流上板，检验校准模型对二者的预测（清单 B08、B09；另建议同期在 INT8 比特流上完成 B12）。
  \item \textbf{M3　向量通路再配比}：在校准模型中完成“链数 $\times$ 向量 lane 数 $\times$ 时钟”的设计空间扫描，选出候选点综合，上板验证。
  \item \textbf{M4　数值方案升级}：完成逐层敏感度分析；BFP8 链上板并验证精度（清单 B11）。
  \item \textbf{M5　片上驻留与主机开销}：前缀 KV 驻留原型；主机调用路径优化。
  \item \textbf{M6　闭环评测}：在仿真环境中完成闭环评测，为真实机器人演示做准备。
\end{enumerate}
每个里程碑都以第~\ref{sec:verif} 节的验证链为准入条件：任何新设计点在进入下一阶段之前，必须在整 chunk 上与期望内存位真，并给出 46 帧精度。

%======================================================================
\section{结论}\label{sec:concl}
%======================================================================

本报告说明了一个在单片 Speedster7t AC7t1500 上端到端执行完整 \pizero{} chunk 的设计方案与验证体系。架构以 35 个自主节点为基础，以列并行 MLP72 级联链承担矩阵运算、以独立向量节点承担非线性与量化运算、以 GDDR6 标志字实现分布式同步，并由 chunk 生成器把整个 chunk 编译为各节点的程序与期望内存。验证体系以五级验证链为骨架，把精度问题与正确性问题分离，并以逐拍核对支撑积极的调度优化。在最佳 INT8 设计点上，整 chunk 时延 1.221~s、关节最大误差 $0.665^\circ$，所有列出设计点均与期望内存位真；与 Jetson AGX Thor 的 125.3~ms 相比约慢 9.7 倍 \tD，差距主要来自向量通路与串行权重装载，以及布线对阵列规模的限制。演示方案以录制数据回放与逐拍核对为核心，突出完整性、精度与确定性。后续工作应优先投入向量通路的结构性扩容与 epilogue 融合，并补齐闭环精度、能耗与时延分布等关键测量。


从设计评审的角度，可以对本项目作出以下总体判断。

\textbf{功能与正确性：已达成。} 完整的 \pizero{} chunk 在单片 FPGA 上端到端运行，主机只负责输入输出；所有列出设计点在完整阵列上与期望内存逐位一致 \tM。五级验证链把精度与正确性干净地分离，使 INT4 这类“硬件正确但精度不可用”的结论可以被明确地得出，也使调度可以被积极优化而不牺牲可信度。

\textbf{精度：在 INT8 上达成，低位宽尚未解决。} W8A8 精确路径在 46 帧上的误差为亚度量级 \tM，可以作为后续一切优化的精度基线；均匀 INT4 不可用，块浮点是最有希望的替代方向，但仍需上板确认。

\textbf{性能：瓶颈明确，差距仍大。} 与 Thor 参照相比约慢 9.7 倍 \tD；校准模型给出的瓶颈（向量通路与串行权重装载）与各项实测结果的方向一致。GEMM 侧的优化（低位宽、可重构）在当前瓶颈下收益受限，这一判断已经被多组数据支持。

\textbf{实现：布线是关键约束。} 阵列规模受布线而非乘法器限制；fabric-light 链证明了布线代价可以大幅降低，但其时钟需要恢复。实现层面的下一步工作是时序收敛与“矩阵/向量”资源再配比。

\textbf{演示：具备条件。} 以最佳 INT8 设计点为基础、以录制数据回放与逐拍核对为核心的演示方案不依赖尚未完成的工作，可以在现有构建产物上实施；演示中的批量回放还将补齐时延分布这一验证缺口。

综上，本项目已经完成了“能否把一个完整的数十亿参数 VLA 正确地放上单片 FPGA”这一问题的回答，并建立了一套能够支撑后续快速迭代的设计与验证基础设施；下一阶段的问题是“能否把它做得足够快”，而这一问题的答案，首先取决于向量通路与片上存储的重新设计。

\appendix
%======================================================================
\section{评审检查清单}\label{app:checklist}
%======================================================================

为便于设计评审，本附录把报告中的关键检查项汇总为清单。

\textbf{架构与微结构。}
\begin{itemize}
  \item 链的长度与位置是否与 MLP72 的列结构对齐；硬级联是否只用于同列相邻单元之间。
  \item 每类节点的数量配比是否有校准模型的依据；向量通路的吞吐是否与矩阵阵列匹配。
  \item 同步标志的放置是否避开大流量的外存通道。
  \item 跨时钟域的每一条通路是否有异步 FIFO 或同步器，是否对频率比无假设。
\end{itemize}

\textbf{数值。}
\begin{itemize}
  \item 位真模型的整数宽度、舍入与饱和是否与 RTL 完全一致。
  \item epilogue 与向量节点的同名运算是否由同一份参数化描述生成。
  \item K 分段的合并顺序是否在位真模型与编译器中一致。
  \item 查表内容是否由位真模型生成并纳入版本管理。
\end{itemize}

\textbf{编译器。}
\begin{itemize}
  \item 所有跨节点的数据依赖是否都有 \texttt{WAIT}/\texttt{POST} 保护（静态检查）。
  \item 张量的地址区间在其生命周期内是否互不重叠（静态检查）。
  \item 期望内存是否与独立的位真模型结果一致。
\end{itemize}

\textbf{验证。}
\begin{itemize}
  \item 新设计点是否在完整阵列、整 chunk、全部 46 帧上逐位一致。
  \item 是否给出了 46 帧的精度统计，以及与上一设计点的比较。
  \item 是否记录了时延与校准模型预测的偏差。
  \item 是否定期进行错误注入测试。
\end{itemize}

\textbf{实现。}
\begin{itemize}
  \item 时序是否在所有三个时钟域上收敛；关键路径是否已知。
  \item 资源占用是否记录在按节点类型的资源预算表中。
  \item 构建产物是否完整可追溯。
\end{itemize}

\textbf{演示。}
\begin{itemize}
  \item 演示所用构建产物是否已完成全部 46 帧的回放与核对。
  \item 界面上的每个数字是否标明实测或模型预测。
  \item 是否准备了后备方案（备用比特流、仿真录屏）。
\end{itemize}


\bibliographystyle{unsrtnat}
\bibliography{vla_fpga_survey}

\end{document}
